% !TeX program = xelatex
\documentclass[12pt]{nextbook}

\UseTemplateSet{
  globalfonts = {libertinus,tibetan},
  features = {hyperlinks}
}

\newcommand{\ReaderVocabBaselineskip}{15pt}
\ExplSyntaxOn
\NewDocumentCommand{\ReaderTibetanBreaks}{m}
  {
    \tl_set:Nn \l_tmpa_tl {#1}
    \regex_replace_all:nnN { ་|༌ } { \0\c{allowbreak} } \l_tmpa_tl
    \regex_replace_all:nnN { \c{allowbreak}([།༎༏༐༑༔]) } { \1 } \l_tmpa_tl
    \tl_use:N \l_tmpa_tl
  }
\ExplSyntaxOff
\newcommand{\ReaderColumnSep}{8mm}
\newcommand{\ReaderVocabText}[1]{\allowbreak\textbf{\ReaderTibetanBreaks{#1}}}
\newcommand{\ReaderVocabTerm}[1]{\textbf{\ReaderTibetanBreaks{#1}}}
\newcommand{\ReaderVocabSeparator}{}
\usepackage{paracol}
\usepackage{xcolor}
\usepackage{eso-pic}
\usepackage{needspace}

\vbadness=10000
\hbadness=10000
\hfuzz=3pt
\emergencystretch=1.5em

\makeatletter
\AtBeginDocument{%
  \@ifpackageloaded{hyperref}{%
    \pdfstringdefDisableCommands{%
      \def\ruby#1#2{#1}%
      \def\ReaderVocabText#1{#1}%
      \def\ReaderVocabTerm#1{#1}%
    }%
  }{}%
}
\makeatother

\providecommand{\ReaderColumnSep}{6mm}
\providecommand{\ReaderColumnRatio}{0.75}
\providecommand{\ReaderVocabBaselineskip}{20pt}
\providecommand{\ReaderVocabText}[1]{\textbf{#1}}
\providecommand{\ReaderVocabTerm}[1]{\textbf{#1}}
\providecommand{\ReaderVocabSeparator}{:}
\providecommand{\ReaderVocabItem}[2]{\noindent\ReaderVocabTerm{#1}\ReaderVocabSeparator\ #2\par}

\setlength{\columnsep}{\ReaderColumnSep}
\setlength{\columnseprule}{0pt}
\columnratio{\ReaderColumnRatio}
\swapcolumninevenpages

\newcounter{readersection}
\newwrite\vocabfile
\newif\iffirstreaderlabelwritten
\newlength{\readercolumnswidth}
\newlength{\readertitlegutter}
\setlength{\readercolumnswidth}{\dimexpr\textwidth-\columnsep\relax}
\setlength{\readertitlegutter}{.75\columnsep}

\makeatletter
\newcommand{\readerseparatorline}{%
  \if@mainmatter
    \AtTextLowerLeft{%
      \ifodd\value{page}%
        \put(\LenToUnit{\dimexpr .75\readercolumnswidth+.5\columnsep\relax},0){%
          \color{black!35}\rule{0.3pt}{\textheight}%
        }%
      \else
        \put(\LenToUnit{\dimexpr .25\readercolumnswidth+.5\columnsep\relax},0){%
          \color{black!35}\rule{0.3pt}{\textheight}%
        }%
      \fi
    }%
  \fi
}
\makeatother

\AddToShipoutPictureBG{\readerseparatorline}

\newcommand{\currentvocabfile}{\jobname-vocab-\thereadersection.voc}
\newcommand{\beginvocabwrite}{\immediate\openout\vocabfile=\currentvocabfile}
\newcommand{\finishvocabwrite}{\immediate\closeout\vocabfile}
\newcommand{\vocabitem}[2]{\ReaderVocabItem{#1}{#2}}
\newcommand{\printvocablist}{%
  \footnotesize
  \setlength{\parindent}{0pt}%
  \setlength{\parskip}{0pt}%
  \setlength{\baselineskip}{\ReaderVocabBaselineskip}%
  \raggedright
  \input{\currentvocabfile}%
}

\newcommand{\readersectiontitle}[1]{%
  \par
  \begingroup
  \ifodd\value{page}%
    \leftskip=\readertitlegutter%
    \rightskip=0.25\textwidth plus 1fil%
  \else
    \leftskip=\dimexpr .25\readercolumnswidth+.5\columnsep+\readertitlegutter\relax%
    \rightskip=0pt plus 1fil%
  \fi
  \section{#1}%
  \par
  \endgroup
}

\NewDocumentCommand{\voc}{omm}{%
  \ReaderVocabText{#2}%
  \IfNoValueTF{#1}{%
    \immediate\write\vocabfile{\string\vocabitem{#2}{#3}}%
  }{%
    \immediate\write\vocabfile{\string\vocabitem{#1}{#3}}%
  }%
}

\newenvironment{readersection}[1]{%
  \stepcounter{readersection}%
  \beginvocabwrite
  \Needspace{10\baselineskip}%
  \iffirstreaderlabelwritten\else
    \label{reader:firstpage}%
    \global\firstreaderlabelwrittentrue
  \fi
  \readersectiontitle{#1}%
  \begin{paracol}{2}
}{%
  \finishvocabwrite
  \switchcolumn
  \printvocablist
  \end{paracol}
}


\title{བོད་སྐད་རྙིང་པའི་གཞི་རིམ་ཀློག་དེབ།}
\author{ནམ་མཁའ་}
\date{\today}

\begin{document}
\frontmatter
\maketitle
\tableofcontents
\mainmatter
\chapter{གཞི་རིམ་ཀློག་ཚན།}
\begin{readersection}{རྒྱལ་པོ་དེ}

\voc{རྒྱལ་པོ་}{king} \voc{ཆེན་པོ་}{great} \voc{དེ་}{that}།
\voc{རྒྱལ་བུ་}{prince} \voc{ཐ་ཆུངས་}{youngest} དེ།
\voc{སྲས་}{son} ཐ་ཆུངས་དེ།
\voc{བུ་མོ་}{daughter} \voc{ངན་པ་}{bad} ཞིག།
བུ་མོ་དེ།
རྒྱལ་བུ་དེ།
སྲས་དེ།
རྒྱལ་པོ་དེ།

\voc{ཆོས་}{teaching} \voc{ཟབ་མོ་}{profound}།
\voc{ཚིག་}{word} \voc{སྙན་པ་}{pleasant}།
ཚིག་ \voc{འཇམ་པ་}{gentle}།
ཚིག་ངན་པ།
ཆོས་ཟབ་མོ་དེ།
ཚིག་སྙན་པ་དེ།
ཚིག་འཇམ་པ་དེ།
\voc{ཉེས་པ་}{evil} ཆེན་པོ་དེ།
ཉེས་པ་ངན་པ་ \voc{དུ་མ་}{many}།

\voc{བྲག་}{rock} \voc{རོག་པོ་}{black} ཞིག།
བྲག་ཆེན་པོ་དེ།
བྲག་རོག་པོ་དེ།
བྲག་ཆེན་པོ་དུ་མ།
\voc{ཚལ་བ་}{splinter} \voc{རྣོན་པོ་}{sharp} དུ་མ།
ཚལ་བ་ \voc{རློན་པ་}{fresh} \voc{གསུམ་}{three}།
ཚལ་བ་རྣོན་པོ་དེ།
ཚལ་བ་རློན་པ་དེ།
ཚལ་བ་ངན་པ་དུ་མ།

\voc{ཤ་}{flesh} རློན་པ།
\voc{ཁྲག་}{blood} \voc{དྲོན་མོ་}{warm}།
ཤ་རློན་པ་དེ།
ཁྲག་དྲོན་མོ་དེ།
ཤ་རློན་པ་དུ་མ།
ཁྲག་དྲོན་མོ་དུ་མ།
ཤ་ངན་པ་ཞིག།
ཁྲག་ངན་པ་དེ།

\voc{ཞག་}{day} དུ་མ།
\voc{ལོ་}{year} གསུམ།
\voc{ཡུན་}{time} \voc{རིང་པོ་}{long} དེ།
ཞག་གསུམ།
ལོ་དུ་མ།
ཡུན་རིང་པོ་དུ་མ།
ཞག་རིང་པོ་དེ།
ལོ་རིང་པོ་དེ།
ཡུན་ཆེན་པོ་དེ།

\voc{ཕུག་རོན་}{dove} གསུམ་ཞིག།
ཕུག་རོན་དུ་མ།
ཕུག་རོན་དེ།
ཕུག་རོན་གསུམ་དེ།
ཕུག་རོན་ཆེན་པོ་ཞིག།
ཕུག་རོན་ངན་པ་ཞིག།
ཕུག་རོན་རློན་པ་མིན་པ།

རྒྱལ་པོ་ཆེན་པོ་དེ།
རྒྱལ་བུ་ཐ་ཆུངས་དེ།
སྲས་ཐ་ཆུངས་དེ།
བུ་མོ་ངན་པ་དེ།
ཆོས་ཟབ་མོ་དེ།
ཚིག་སྙན་པ་དེ།
བྲག་རོག་པོ་དེ།
ཚལ་བ་རྣོན་པོ་དེ།
ཤ་རློན་པ་དེ།
ཁྲག་དྲོན་མོ་དེ།
ཡུན་རིང་པོ་དེ།
ཕུག་རོན་གསུམ་དེ།

\end{readersection}

\begin{readersection}{སེམས་ཅན་ཆེན་པོ}

\voc{བདག་}{I} སེམས་ཅན་ཐ་ཆུངས་མིན།
\voc{ང་}{I} \voc{སེམས་ཅན་ཆེན་པོ་}{Sems-chan-chen-po} \voc{ཡིན་}{to be}།
སེམས་ཅན་ཆེན་པོ་དེ་ རྒྱལ་བུ་མིན།
དེ་ \voc{ཐུགས་རྗེ་ཅན་}{compassionate} ཡིན།
དེ་ \voc{མགོན་པོ་}{protector} \voc{བཟང་པོ་}{excellent} ཡིན།
དེ་ མགོན་པོ་ངན་པ་མིན།

རྒྱལ་པོ་ཆེན་པོ་དེ་ \voc{དྲག་པ་}{noble} ཡིན།
རྒྱལ་པོ་དེའི་ \voc{འཁོར་}{attendant} \voc{མང་པོ་}{many} ཡིན།
འཁོར་དེ་དག་ \voc{ཀུན་}{all} བཟང་པོ་མིན།
འཁོར་དེ་དག་ལ་ ཚིག་སྙན་པ་ \voc{མང་དུ་}{much}།
ཚིག་ངན་པ་ཡང་མང་དུ།

\voc{ཁྱིམ་མཚེས་}{neighbour} དེ་ \voc{དགེ་བསྙེན་}{layman} ཡིན།
དགེ་བསྙེན་དེ་ རྒྱལ་པོ་མིན།
དགེ་བསྙེན་དེ་ འཁོར་མིན།
དེ་ སེམས་ཅན་ཐུགས་རྗེ་ཅན་ཡིན།
དེའི་ \voc{བརྡ་}{sign} དེ་ ཚིག་སྙན་པ་མིན།
བརྡ་དེ་ ཚིག་འཇམ་པ་ཡིན།

\voc{སག་མོ་}{tigress} དེ་ \voc{ལྟོགས་པ་}{hungry} ཡིན།
སག་མོ་དེ་ སེམས་ཅན་ངན་པ་མིན།
\voc{སྤེའུ་ཕྲུག་}{monkey-child} དེ་ སེམས་ཅན་ཆུང་ངུ་ཡིན།
སྤེའུ་ཕྲུག་དེ་ སག་མོ་མིན།
\voc{སྲོག་ཆགས་}{animal} དེ་དག་ སེམས་ཅན་ཡིན།
སྲོག་ཆགས་ཀུན་ སེམས་ཅན་ཡིན།

རྒྱལ་པོ་དེའི་ \voc{སོན་མོ་}{feast} ཆེན་པོ་ཡིན།
སོན་མོ་དེ་ སོན་མོ་ངན་པ་མིན།
སོན་མོ་དེ་ལ་ \voc{ལྟོ་གོས་}{food and clothing} མང་པོ།
ལྟོ་གོས་དེ་ སེམས་ཅན་ལྟོགས་པ་ཀུན་གྱི་ཡིན།
དེ་ ཆོས་ཟབ་མོ་མིན།
འོན་ཀྱང་དེ་ སེམས་ཅན་ལ་ བཟང་པོ་ཡིན།

སེམས་ཅན་ \voc{ཁྲི་}{ten thousand} དང་ སྲོག་ཆགས་ \voc{སོང་}{thousand} ཡིན།
འཁོར་ \voc{ལྔ་བརྒྱ་}{five hundred} ཡིན།
སག་མོ་གསུམ་ཡིན།
སྤེའུ་ཕྲུག་དུ་མ་ཡིན།
དེ་དག་ཀུན་ སེམས་ཅན་ཡིན།
སེམས་ཅན་ཆེན་པོ་དེ་ ཀུན་གྱི་མགོན་པོ་ཡིན།

\end{readersection}

\begin{readersection}{ཁ་བ་ཅན་གྱི་ཡུལ}

\voc{ཁ་བ་ཅན་}{Tibet} གྱི་ \voc{ཡུལ་}{country} དེ་ན་ \voc{གྲོང་ཁྱེར་}{village} ཆེན་པོ་ཞིག་ \voc{ཡོད་}{to exist}།
གྲོང་ཁྱེར་དེ་ན་ རྒྱལ་པོ་དྲག་པ་ཞིག་ཡོད།
རྒྱལ་པོ་དེ་ལ་ རྒྱལ་བུ་གསུམ་ཡོད།
རྒྱལ་བུ་ཐ་ཆུངས་དེ་ སེམས་ཅན་ཆེན་པོ་མིན།
འོན་ཀྱང་དེ་ སེམས་ཅན་ཐུགས་རྗེ་ཅན་ཡིན་ནོ།

གྲོང་ཁྱེར་དེ་ན་ \voc{རྒན་མོ་}{old woman} ཞིག་ཡོད།
རྒན་མོ་དེ་ལ་ \voc{བུ་}{son} གཉིས་དང་ \voc{ཕྲུ་གུ་}{child} གསུམ་ཡོད།
བུ་དེའི་ \voc{སྐད་}{voice} སྙན་པ་ཡིན།
ཕྲུ་གུ་དེ་དག་གི་སྐད་ འཇམ་པ་མིན།
འོན་ཀྱང་ཕྲུ་གུ་དེ་དག་ ངན་པ་མིན་ནོ།

\voc{ཁྱོད་}{you} གྲོང་ཁྱེར་དེ་ན་ཡོད།
\voc{བདག་ཅག་}{we} ཡུལ་དེ་ན་ཡོད།
བདག་ཅག་ལ་ ལྟོ་གོས་མང་པོ་ཡོད།
ཁྱོད་ལ་ ཚིག་སྙན་པ་དུ་མ་ཡོད།
བདག་ཅག་གི་ཚིག་ ཚིག་ངན་པ་མིན་ནོ།

གྲོང་ཁྱེར་དེ་ལ་ \voc{ཆུ་}{water} མང་པོ་ཡོད།
ཆུ་དེ་ ཁྲག་དྲོན་མོ་མིན།
ཆུ་དེ་ ཆུ་འཇམ་པ་ཡིན།
\voc{རྒྱ་མཚོ་}{ocean} ཆེན་པོ་ལ་ ཆུ་ \voc{དཔག་མེད་}{countless} ཡོད།
ཡུལ་དེ་ལ་ རྒྱ་མཚོ་ \voc{མེད་}{not exist}།
འོན་ཀྱང་གྲོང་ཁྱེར་དེ་ལ་ ཆུ་ཡོད་དོ།

\voc{བྲམ་ཟེ་}{Brahmin} དེ་ དགེ་བསྙེན་མིན།
བྲམ་ཟེ་དེ་ལ་ \voc{ཆུང་མ་}{wife} བཟང་པོ་ཡོད།
ཆུང་མ་དེ་ \voc{བུད་མེད་}{woman} བཟང་པོ་ཡིན།
བུད་མེད་དེའི་ \voc{ཆས་}{garment} དྲོན་མོ་ཡིན།
བྲམ་ཟེ་དེའི་ཆུང་མ་ལ་ ཚིག་འཇམ་པ་ཡོད།
བྲམ་ཟེ་དེ་ལ་ \voc{སྡོམ་པ་}{vow} ཟབ་མོ་ཡོད་དོ།

\voc{འགྲོ་བ་}{being} ཀུན་ལ་ \voc{ཚེ་སྲོག་}{life} ཡོད།
ཚེ་སྲོག་མེད་པ་དེ་ སེམས་ཅན་མིན།
\voc{ཚེ་}{time} རིང་པོ་དེ་ལ་ ཉེས་པ་མང་པོ་མེད།
རྒན་མོ་དེ་ལ་ ཚེ་རིང་པོ་ཡོད།
ཕྲུ་གུ་དེ་དག་ལ་ ལོ་གསུམ་མེད།
བུ་དེ་ལ་ ལོ་དུ་མ་ཡོད།

\voc{སྤྱན་རས་གཟིགས་}{Avalokiteśvara} ལ་ ཐུགས་རྗེ་ཟབ་མོ་ཡོད།
སྤྱན་རས་གཟིགས་ འགྲོ་བ་ཀུན་གྱི་ མགོན་པོ་ཡིན།
སེམས་ཅན་ལྟོགས་པ་ཀུན་ལ་ མགོན་པོ་ཡོད།
སག་མོ་དེ་ལ་ ཤ་རློན་པ་མེད།
སྤེའུ་ཕྲུག་དེ་ལ་ ལྟོ་གོས་མང་པོ་མེད།
འོན་ཀྱང་སེམས་ཅན་ཆེན་པོ་དེ་ལ་ ཐུགས་རྗེ་ཡོད་དོ།

\voc{རི་རབ་}{Sumeru} ལ་ \voc{གཏེར་}{treasury} དཔག་མེད་ཡོད།
རི་རབ་དེ་ལ་ \voc{སྤེའུ་}{monkey} མང་པོ་ཡོད།
སྤེའུ་དེ་དག་གི་ སྐད་ སྙན་པ་མིན།
\voc{ཁོང་སེང་}{cavern} དེ་ན་ \voc{འབྱུང་པོ་}{demon} ཡོད།
ཁོང་སེང་གཞན་དེ་ན་ \voc{སྲིན་པོ་}{demon} མེད།
འབྱུང་པོ་དེ་ སེམས་ཅན་བཟང་པོ་མིན་ནོ།

\voc{ཨ་ཁུ་}{uncle} དེ་ བྲམ་ཟེ་ཡིན།
ཨ་ཁུ་དེའི་ བུ་ \voc{ཕྱིན་ཏེ་ལོ་ཤུ་ཤ་}{Phyin-te-lo-shu-sha} ཡིན།
ཕྱིན་ཏེ་ལོ་ཤུ་ཤ་དེ་ལ་ ཆས་བཟང་པོ་ཡོད།
\voc{བྱ་}{bird} དེ་ ཕུག་རོན་ཡིན།
ཕུག་རོན་དེ་ལ་ སྐད་སྙན་པ་ཡོད།
ཕུག་རོན་གསུམ་དེ་ གྲོང་ཁྱེར་དེ་ན་ཡོད་དོ།

\end{readersection}

\begin{readersection}{ཚལ་དུ་སོང}

ཁ་བ་ཅན་གྱི་ཡུལ་ན་ \voc{གནས་}{place} ཆེན་པོ་ཞིག་ཡོད།
གནས་དེ་ན་ \voc{ཚལ་}{forest} དང་ བྲག་ཆེན་པོ་ཡོད།
བྲག་དེ་ལ་ \voc{བྲག་སྲིན་མོ་}{rock-ogress} ཞིག་ \voc[སྐྱེ་]{སྐྱེས་}{to be born}།
\voc{དེར་}{there} མི་རིགས་ཀྱི་ \voc{ཡབ་མེས་}{ancestors} \voc[འབྱུང་]{བྱུང་}{to occur}།
མི་རིགས་དེ་ ངན་པ་མིན།
མི་རིགས་དེ་ལ་ ཚིག་སྙན་པ་དང་ ཆོས་ཟབ་མོ་ཡོད།

\voc{མི་}{man} \voc{རིགས་}{race} དེ་ལ་ \voc{བུ་ཕོ་}{son} \voc{གཉིས་}{two} ཡོད།
བུ་ཕོ་གཉིས་ཀྱི་ཨ་ཁུ་ དགེ་བསྙེན་ཡིན།
ཨ་ཁུ་དེ་ མགོན་པོ་མིན།
འོན་ཀྱང་དེ་ ཐུགས་རྗེ་ཅན་ཡིན།
དེ་ལ་ ལྟོ་གོས་མང་པོ་ཡོད།
དེ་ལ་ \voc{ནོར་}{wealth} དུ་མ་ཡོད།

བུ་ཕོ་གཉིས་ ཞག་གསུམ་གྱི་ཡུན་ལ་ གྲོང་ཁྱེར་དེ་ན་ཡོད།
\voc{དེ་ནས་}{then} བུ་ཕོ་གཉིས་ \voc[ལང་]{ལངས་}{to rise}།
བུ་ཕོ་ཆེན་པོ་ ཡར་ལངས།
བུ་ཕོ་ཐ་ཆུངས་ ཕྱིར་ལངས།
བུ་ཕོ་གཉིས་ ཚལ་དུ་ \voc[འགྲོ་]{སོང་}{to go}།
ཨ་ཁུ་དེ་ཡང་ ཚལ་དུ་ \voc[འདོང་]{དོང་}{to go}།

ཚལ་དེ་ན་ \voc{ལོ་ཐོག་}{harvest} \voc{གང་བ་}{full} ཡོད།
ཚལ་དེ་ན་ \voc{ཤིང་ཐོག་}{fruit} དུ་མ་ཡོད།
ཤིང་ཐོག་དེ་ རློན་པ་དང་ དྲོན་མོ་མིན།
ཆུ་དེ་ན་ ཕུག་རོན་གསུམ་ཡོད།
བྱ་དེ་དག་ སྐད་སྙན་པའི་སྲོག་ཆགས་ཡིན།
སག་མོ་ལྟོགས་པ་དེ་ ཚལ་དེ་ན་ཡོད།

བུ་ཕོ་གཉིས་ཀྱི་ ལྟོ་གོས་ \voc[འཛད་]{ཟད་}{to be consumed}།
ཤིང་ཐོག་མང་པོ་ཡང་ཟད།
དེ་ནས་ བུ་ཕོ་གཉིས་ལ་ ལྟོ་གོས་ \voc{མེད་}{not have}།
ཨ་ཁུ་དེ་ལ་ ནོར་མང་པོ་མེད།
འོན་ཀྱང་ཨ་ཁུ་དེ་ལ་ ཆོས་ཟབ་མོ་ཡོད།

བྲག་ཆེན་པོ་དེ་ལ་ \voc{གོས་}{clothes} \voc{ངུར་སྨྲིག་}{reddish yellow} ཞིག་ \voc{འདུག་}{to lie}།
གོས་དེའི་ནང་ན་ \voc{རུས་བུ་}{small bone} ཞིག་འདུག།
རུས་བུ་དེ་ སྲོག་ཆགས་ཀྱི་མིན།
དེ་ \voc{སྲིན་མོ་}{demoness} དེའི་རུས་བུ་ཡིན།
སྲིན་མོ་དེ་ \voc[ཚེ་འཕོ་]{ཚེ་འཕོས་}{to die}།
སྲིན་པོ་དེ་ \voc{ཕྱིར་}{back} དོང།

\voc{འཕགས་པ་}{Noble One} དེ་ \voc{བཞེངས་}{to rise}།
དེ་ནས་ འཕགས་པ་དེ་ གནས་དེར་ \voc[འོང་]{འོངས་}{to come}།
\voc{ལྷ་}{deity} གཉིས་ཀྱང་ གནས་དེར་འོངས།
ལྷ་གཉིས་ལ་ ཐུགས་རྗེ་ཟབ་མོ་ཡོད།
གནས་དེ་ བཟང་པོར་ \voc[འགྱུར་]{གྱུར་}{to become}།
ཚལ་དེ་ གནས་བཟང་པོར་གྱུར།

དེ་ནས་ བུ་ཕོ་གཉིས་ གྲོང་ཁྱེར་དུ་ ཕྱིར་འོངས།
ཨ་ཁུ་དེ་ཡང་ ཕྱིར་འོངས།
རྒན་མོ་དེ་ལ་ བུ་ཕོ་གཉིས་ཀྱི་ཚིག་སྙན་པ་བྱུང།
བུ་ཕོ་གཉིས་ཀྱི་ ཆས་ ངན་པར་གྱུར།
འོན་ཀྱང་དེ་དག་གི་ སེམས་ བཟང་པོར་གྱུར།

གནས་དེ་ན་ ཉེས་པ་ཆེན་པོ་མེད།
ཚལ་དེ་ན་ འབྱུང་པོ་མེད།
ཁོང་སེང་དེ་ན་ སྲིན་པོ་མེད།
ལྷ་གཉིས་ \voc{ཡར་}{up} སོང།
འཕགས་པ་དེ་ གནས་དེར་འདུག།
མི་རིགས་དེ་ཡང་ གནས་དེར་སྐྱེས་ཤིང་ བཟང་པོར་གྱུར།

\end{readersection}

\begin{readersection}{དབྱིག་པ་ཅན་}
ཁ་བ་ཅན་གྱི་ཡུལ་དུ་གྲོང་ཁྱེར་ཆེན་པོ་གཅིག་ཡོད་དོ། གྲོང་ཁྱེར་དེའི་\voc{སྒོ་}{door}ན་ཁྱིམ་མཚེས་མང་པོ་\voc[སྡོད་]{བསྡད་}{to stay}དོ། དེར་\voc{དབྱིག་པ་ཅན་}{Dbyig-pa-čan}ཡོད། དབྱིག་པ་ཅན་དེར་\voc{དོ་སྨད་}{family}དང་\voc{ཕ་མ་}{parents}དང་\voc{ཕ་མིང་}{father’s brothers}དང་\voc{ཨ་ནེ་}{aunt}དང་\voc{ཞང་པོ་}{uncle}ཡོད་དོ། དེ་ཀུན་ལ་\voc{བུ་ཚ་}{offspring}མང་པོ་ཡོད། དབྱིག་པ་ཅན་གྱི་\voc{ཁྱོ་}{husband}ནི་\voc{ཐ་ག་པ་}{weaver}ཡིན། ཐ་ག་པ་དེའི་\voc{སྐྲ་}{hair}རིང་པོ་དང་\voc{ཁ་སྤུ་}{beard}རོག་པོ་ཡོད། དེའི་\voc{རྐང་པ་}{leg}\voc{གཅིག་}{one}ནི་\voc[འཆག་]{ཆག་}{to break}ནས་\voc{འདམ་}{swamp}དུ་འདུག་གོ།

དེ་ནས་རྒྱལ་པོས་རྒྱལ་བུ་གཉིས་དང་འཁོར་མང་པོ་དང་ལྷན་ཅིག་གྲོང་ཁྱེར་གྱི་སྒོར་སོང་། རྒྱལ་པོས་དབྱིག་པ་ཅན་གྱི་ཁྱིམ་ལ་\voc[ལྟ་]{བལྟས་}{to look}། དབྱིག་པ་ཅན་གྱིས་རྒྱལ་པོ་དང་རྒྱལ་བུ་ལ་གོས་ངུར་སྨྲིག་དང་ལྟོ་གོས་མང་པོ་\voc[འབྱིན་]{ཕྱུང་}{to take out}ངོ་། ཐ་ག་པས་གོས་དེ་རྒྱལ་བུ་ལ་\voc[སྟོན་]{བསྟན་}{to teach}། རྒྱལ་བུས་གོས་ལ་བལྟས་ནས་གོས་འདི་བཟང་པོ་ཡིན་ནོ་\voc{སྙམ་}{thought}དུ་\voc[སེམས་]{བསམས་}{to think}སོ། དབྱིག་པ་ཅན་གྱིས་ཁྱིམ་གྱི་ནོར་དང་གོས་\voc[འདྲིལ་]{དྲིལ་}{to unite}ནས་\voc{གཅིག་པོ་}{only}ཁྱོ་ལ་མ་བྱིན། དབྱིག་པ་ཅན་གྱིས་ནོར་ཀུན་སེམས་ཅན་ལྟོགས་པ་ལ་བྱིན་ནོ།

གྲོང་ཁྱེར་དེའི་ཕྱིར་ན་\voc{རྟ་}{horse}\voc{གླང་}{ox}དང་ལྷན་ཅིག་ཆུ་ལ་སོང་། རྟ་དེ་ནི་\voc{རྟག་ཏུ་}{always}སྒོ་དེར་འོང་། གླང་དེ་འདམ་དུ་\voc[འཚུད་]{ཚུད་}{to enter}དོ། ཐ་ག་པས་རྟའི་ལུས་ལ་བལྟས། དབྱིག་པ་ཅན་གྱིས་གླང་གི་\voc{ལུས་}{body}ལ་\voc[འཛིན་]{བཟུང་}{to seize}ནས་ཕྱིར་ཕྱུང་ངོ་། དེར་སྤེའུ་ཕྲུག་ཅིག་གིས་\voc{ཤིང་}{wood}ལ་འཛེགས་ནས་བྱི་བར་གྱུར། བྱ་གཅིག་ཀྱང་ནམ་མཁའ་ནས་\voc[འབྱི་]{བྱི་}{to fall off}སྟེ་ཆུར་ཚུད་དོ། རྒྱལ་བུས་བྱ་དེ་ཕྱུང་ནས་གནས་བཟང་པོར་བཞག་གོ།

དེ་ནས་\voc{བཅོམ་ལྡན་འདས་}{Buddha}དེར་འཁོར་མང་པོ་དང་ལྷན་ཅིག་བྱུང་ངོ་། བཅོམ་ལྡན་འདས་ཀྱིས་\voc{འཁོར་བ་}{existence}དང་\voc{སྐྱེ་བ་}{birth}དང་\voc{འཇིག་པ་}{decay}རྣམས་མཁྱེན་ནས་འཁོར་ལ་ཆོས་བསྟན། དེའི་\voc{ཐུགས་}{mind}ནི་ཐུགས་རྗེ་ཅན་ཡིན། བཅོམ་ལྡན་འདས་ཀྱིས་དབྱིག་པ་ཅན་ལ་སྐྱེ་བ་དང་འཁོར་བ་དང་འཇིག་པ་བསྟན་ཏོ། དབྱིག་པ་ཅན་གྱིས་ཆོས་དེ་\voc{མཁྱེན་}{to recognise}ནས་\voc{ཡིད་}{mind}བཟང་པོར་གྱུར། དབྱིག་པ་ཅན་གྱིས་\voc{བསམ་པ་}{thought}བཟང་པོ་བསམས་ཏེ། སེམས་ཅན་\voc{ཀུན་}{all}གྱི་\voc{རང་}{itself}གི་ལུས་དང་ཡིད་སྲོག་ཆགས་དང་འདྲའོ་སྙམ་མོ།

ཡུལ་དེར་\voc{མཛེས་སེ་}{Mdzesse}ཡོད། མཛེས་སེ་དེས་དབྱིག་པ་ཅན་ལ་ཚིག་ངན་པ་\voc[སྤྱོ་]{སྤྱོས་}{to scold}སོ། དབྱིག་པ་ཅན་གྱིས་མཛེས་སེ་ལ་མ་སྤྱོས། དབྱིག་པ་ཅན་གྱིས་\voc{གཞན་}{other}གྱི་ཉེས་པ་ལ་བལྟས་ཀྱང་རང་གི་སེམས་བཟང་པོར་བཟུང་ངོ་། དེ་ནས་མཛེས་སེ་ཡང་དབྱིག་པ་ཅན་དང་ལྷན་ཅིག་བཅོམ་ལྡན་འདས་ལ་སོང་། བཅོམ་ལྡན་འདས་ཀྱིས་དེ་གཉིས་ལ་འཁོར་བའི་ཉེས་པ་བསྟན་ཏོ། དེ་གཉིས་ཀྱིས་འཁོར་བ་ན་བདེ་བ་མེད་པ་མཁྱེན་ནས། སེམས་ཅན་ལ་ཆོས་ཟབ་མོ་རྟག་ཏུ་སྟོན་ནོ་སྙམ་མོ།

དེའི་ཚེ་ཡུལ་དེར་སྲིན་པོ་ངན་པ་ཞིག་བྱུང་། སྲིན་པོ་དེས་བྱ་དང་སྲོག་ཆགས་མང་པོ་\voc[གསོད་]{བསད་}{to kill}དོ། རྒྱལ་པོས་སྲིན་པོ་དེ་ལ་བལྟས་ནས་འཇིག་པ་ཆེན་པོ་འབྱུང་སྲིད་དོ་\voc{སྲིད་}{to exist}སྙམ་མོ། དབྱིག་པ་ཅན་གྱིས་སེམས་ཅན་ཀུན་ལ་བསམ་པ་བཟང་པོ་བསམས་ནས། བཅོམ་ལྡན་འདས་དང་འཁོར་ལ་བརྡ་བསྟན། འཁོར་ཀུན་གཅིག་ཏུ་དྲིལ་ནས་སྒོ་ནས་ཕྱིར་སོང་། སྲིན་པོས་འཁོར་ལ་བལྟས་ནས་སེམས་འཇམ་པར་གྱུར། དེ་ནས་སྲིན་པོས་སེམས་ཅན་མ་བསད། རྒྱལ་པོ་དང་དབྱིག་པ་ཅན་དང་མཛེས་སེ་ཀུན་གྱིས་ཆོས་སྟོན་པ་དེ་ལ་བལྟས་ནས་ཐུགས་དགའོ།
\end{readersection}

\begin{readersection}{ཁྱིམ་བདག་གི་ཞལ་ཆེམས་}
ཁ་བ་ཅན་གྱི་ཡུལ་དུ་གྲོང་ཁྱེར་ཆེན་པོ་གཅིག་ཡོད་དོ། གྲོང་ཁྱེར་དེའི་དབུས་ན་ཁྱིམ་ཆེན་པོ་ཞིག་ཡོད། ཁྱིམ་དེའི་\voc{ཁྱིམ་བདག་}{householder}ནི་དགེ་བསྙེན་བཟང་པོ་ཞིག་ཡིན་ནོ། ཁྱིམ་བདག་དེར་ཕ་མ་དང་ཁྱོ་དང་བུ་ཚ་མང་པོ་ཡོད། ཁྱིམ་བདག་གི་\voc{ཡབ་}{father}ནི་རྒན་པོ་ཡིན། ཁྱིམ་བདག་གི་\voc{ཨ་མ་}{mother}ནི་རྒན་མོ་ཡིན། ཁྱིམ་དེ་ན་ནོར་དང་གོས་དང་\voc{དངོས་པོ་}{thing}མང་པོ་ཡོད་ཀྱང་། ཁྱིམ་བདག་དེས་རྟག་ཏུ་སེམས་ཅན་ལྟོགས་པ་ལ་ཟས་དང་གོས་\voc[སྦྱིན་]{བྱིན་}{to give}ནོ།

དེ་ནས་བཅོམ་ལྡན་འདས་འཁོར་མང་པོ་དང་ལྷན་ཅིག་གྲོང་ཁྱེར་དེར་བྱུང་ངོ་། འཁོར་གྱི་\voc{གྲལ་}{row}གྱི་\voc{དབུས་}{middle}ན་\voc{ཀུན་དགའ་བོ་}{Ānanda}བསྡད་དོ། ཀུན་དགའ་བོས་\voc{ཡི་གེ་}{text}ཟབ་མོ་\voc[ཀློག་]{བཀླགས་}{to read}ནས། འཁོར་ཀུན་ལ་ཚིག་སྙན་པ་བསྟན། ཁྱིམ་བདག་གིས་ཡི་གེ་དེ་མཐོང་ནས་ཀུན་དགའ་བོ་ལ་བསམས། འདི་\voc{ད་}{now}ངའི་ཁྱིམ་དུ་\voc{སྐུར་}{to send}ཅིག ངའི་བུ་ཚ་ཀུན་གྱིས་ཆོས་བཀླག་པར་བྱའོ། ཡི་གེ་ཟབ་མོ་འདི་མ་ཉེས་པར་\voc[སྐུར་]{བསྐུར་}{to send}ཅིག་སྙམ་མོ།

ཁྱིམ་བདག་གི་ཁྱིམ་ནས་\voc{ལམ་}{road}རིང་པོ་ཞིག་གྲོང་ཁྱེར་གྱི་སྒོར་འགྲོ། ལམ་དེའི་\voc{བར་}{middle}ན་\voc{ཆུ་བོ་}{river}ཆེན་པོ་ཡོད། ཆུ་བོ་དེ་\voc{ཆུ་}{river}མང་པོ་ལས་\voc{ཁྱབ་}{to comprise}པ་ཡིན་ནོ། ཆུ་བོ་དེའི་ཕྱིར་ན་ཚལ་ཞིག་ཡོད། ཚལ་དེ་ལས་\voc{མེ་ཤིང་}{firewood}མང་པོ་འབྱུང་། ཁྱིམ་བདག་གི་བུ་གཅིག་ཚལ་དེར་སོང་ནས་མེ་ཤིང་བཟུང་ངོ་། ཨ་མས་བུ་ལ་བསམས། ཚལ་དེ་ནས་མེ་ཤིང་ཟུངས་ཤིག འདམ་ནས་མ་འགྲོ། ཆུ་བོ་དེར་མ་ལྟུང་ཞིག་སྙམ་མོ།

བུ་དེས་\voc{སྟེའུ་}{axe}ཞིག་ཁྱིམ་མཚེས་ལས་\voc[བརྙ་]{བརྙས་}{to borrow}ཏེ་ཚལ་དུ་སོང་། དེར་ཤིང་རིང་པོ་ཞིག་མཐོང་ནས་སྟེའུ་ལ་བལྟས། ཤིང་དེ་ལས་\voc{རྡུལ་}{dust}མང་པོ་བྱུང་། རྡུལ་དེས་བུའི་གོས་ལ་\voc{དྲི་མ་}{impurity}བྱུང་། སྟེའུ་དེ་ལ་\voc{སེར་ག་}{crack}ཡང་བྱུང་ངོ་། བུས་སྟེའུ་དེ་ལ་བལྟས་ནས། དེ་ལྟ་བུའི་དངོས་པོ་ཁྱིམ་མཚེས་ལ་ཕྱིར་མ་སྐུར་ཞིག སྟེའུ་གཞན་ཞིག་ཁྱིམ་མཚེས་ལ་བྱིན་ཅིག་སྙམ་མོ།

དེའི་ཚེ་\voc{ལོང་བ་}{blind}རྒན་པོ་ཞིག་ལམ་ནས་འོངས་སོ། ལོང་བ་དེས་ལམ་གྱི་བར་ན་རྡོ་དང་ཤིང་ཡོད་པ་མ་མཐོང་། ཁྱིམ་བདག་གི་བུས་ལོང་བ་དེ་\voc{མཐོང་}{to see}ནས། འདིར་ཤོག ཆུ་བོར་མ་\voc[ལྟུང་]{ལྷུང་}{to fall}ཞིག ལམ་དེ་ནས་མ་འགྲོ། བདག་དང་ལྷན་ཅིག་ཤོག་སྙམ་མོ། ལོང་བ་དེ་བུ་དང་ལྷན་ཅིག་\voc{སླར་}{back}ཁྱིམ་དུ་འོངས་སོ། ལོང་བ་དེའི་\voc{ཡན་ལག་}{limb}ལ་ཉེས་པ་མེད་དོ།

ཁྱིམ་བདག་གི་ཨ་མས་ལོང་བ་ལ་ཟས་དྲོན་མོ་བྱིན་ནས། ཟོ་ཞིག འོན་ཀྱང་ཤ་ངན་པ་མ་\voc[ཟ་]{ཟོ་}{to eat}ཞིག ཟས་འདི་\voc{ཟ་}{to eat}ཞིག་སྙམ་མོ། ལོང་བ་དེས་ཟས་\voc[ཟ་]{བཟས་}{to eat}ནས་དགའོ། ཁྱིམ་བདག་གིས་ལོང་བ་ལ་གོས་\voc{སྡུག་}{pretty}པ་ཞིག་བྱིན། གོས་དེ་ནི་སྡུག་པ་ཡིན་ཡང་། ལོང་བ་དེས་མ་མཐོང་ངོ་། ལོང་བ་དེས། གོས་མི་མཐོང་ཡང་ཁྱོད་ཀྱི་བསམ་པ་བཟང་པོ་མཁྱེན་ནོ་སྙམ་མོ།

དེ་ནས་ཁྱིམ་བདག་གི་ཡབ་ཀྱིས་\voc{ཞལ་ཆེམས་}{testament}བཀླགས་སོ། ཞལ་ཆེམས་དེར། ཁྱིམ་གྱི་ནོར་ཀུན་རང་གཅིག་པོའི་དོན་དུ་མ་འཛིན་ཅིག ནོར་དང་དངོས་པོ་ཀུན་སེམས་ཅན་ལ་སྦྱིན་ཅིག རྡུལ་དང་དྲི་མ་ལས་ཕྱིར་འགྲོ་ཞིག ལམ་ངན་པ་ནས་མ་འགྲོ། ཆོས་ཟབ་མོ་རྟག་ཏུ་ཀློག་ཅིག སེམས་ཅན་ཀུན་ལ་\voc{དེ་ལྟ་བུ་}{such}བའི་བསམ་པ་བཟང་པོ་བསམས་ཤིག་ཅེས་ཡོད་དོ།

ཁྱིམ་བདག་གིས་ཞལ་ཆེམས་དེ་བཀླགས་ནས་ཕ་མ་དང་བུ་ཚ་ཀུན་ལ་བསྟན། ཁྱིམ་ནས་ནོར་མང་པོ་ཕྱུང་སྟེ་སེམས་ཅན་ལ་བྱིན། དེ་ནས་ཁྱིམ་བདག་དང་བུ་ཚ་ཀུན་ཆོས་ཀྱི་ལམ་དུ་\voc[འགྲོ་]{སོང་}{to become}། ཁྱིམ་དེ་ལ་འཁོར་མང་པོ་\voc{འཁོར་}{to roam}ཡང་། དེ་ཀུན་གྱིས་ཆོས་བཀླགས་ནས་སེམས་བཟང་པོར་གྱུར་ཏོ། དངོས་པོ་ཀུན་འཇིག་པར་འགྲོ་བ་ཡིན། ཆོས་ནི་སླར་སླར་ཀློག་པར་བྱའོ།
\end{readersection}

\begin{readersection}{ཉི་མ་ཤར་བ་}
ཁ་བ་ཅན་གྱི་ཡུལ་དུ་གྲོང་ཁྱེར་ཆེན་པོ་གཅིག་ཡོད་དོ། གྲོང་ཁྱེར་དེའི་སྒོ་ན་ཁྱིམ་བདག་བཟང་པོ་ཞིག་བསྡད། ཁྱིམ་བདག་དེའི་ཁྱིམ་ན་\voc{བགྲེས་པ་}{elder}ཞིག་དང་\voc{མ་སྨད་}{mother and children}དང་བུ་ཚ་མང་པོ་ཡོད་དོ། བགྲེས་པ་དེ་ནི་རིགས་བཟང་པོ་ལས་སྐྱེས་པ་ཡིན། དེའི་\voc{རིགས་}{family}ལ་ནོར་མང་པོ་ཡོད་ཀྱང་། བགྲེས་པ་དེས་ནོར་ལ་རྟག་ཏུ་མ་ཆགས། བགྲེས་པ་དེས་ཁྱིམ་གྱི་བུ་ཚ་ལ་ཚིག་བསྟན། འཁོར་བ་ན་སྐྱེ་བ་དང་འཇིག་པ་ཡོད། \voc[སྐྱེ་]{སྐྱེས་}{to be born}པ་ཀུན་ཕྱིར་\voc{འགྲོ་}{to go}བར་འགྱུར། དེ་ལྟ་བུའི་\voc{དོན་}{meaning}མཁྱེན་པར་བྱའོ་སྙམ་མོ།

དེ་ནས་\voc{ཉི་མ་}{sun}\voc[འཆར་]{ཤར་}{to rise}ཏེ། གྲོང་ཁྱེར་དེ་ལ་སེམས་\voc{གསལ་}{clear}པོ་བྱུང་ངོ་། ཉི་མའི་\voc{འོད་ཟེར་}{ray}ཁྱིམ་གྱི་སྒོ་ནས་ཁྱིམ་དུ་ཞུགས། བགྲེས་པ་དེས་འོད་ཟེར་ལ་བལྟས་ནས། ཉི་མ་འཆར་བ་བཟང་པོ་ཡིན། སེམས་གསལ་བ་དེ་ཀོ་ཆོས་ཀྱི་སྒོ་ཡིན་ནོ་སྙམ་མོ། དེར་ཁྱིམ་བདག་གི་བུ་\voc{ལྔ་}{five}དང་ཕྲུ་གུ་\voc{དྲུག་}{six}འདུག བུ་ལྔ་ཀ་དང་ཕྲུ་གུ་དྲུག་ཀས་བགྲེས་པའི་ཚིག་\voc{ཐོས་}{to hear}སོ། བུ་ཐ་ཆུངས་དེས་དེའི་\voc{དེ་སྐད་}{those words}ཐོས་ནས། བདག་ཀོས་ཆོས་ཀློག་པར་བྱའོ་སྙམ་མོ།

ཁྱིམ་དེའི་ཕྱིར་ན་ཚལ་ཆེན་པོ་ཞིག་ཡོད། ཚལ་དེར་\voc{ཤིང་མཁན་}{carpenter}ཞིག་བསྡད། ཤིང་མཁན་དེས་སྟེའུ་བཟུང་ནས་ཤིང་ལ་བལྟས། ཁྱིམ་བདག་གིས་ཤིང་མཁན་དེ་ཁྱིམ་དུ་\voc{ཁྲིད་}{to lead}དེ། ཁྱིམ་གྱི་སྒོ་ལ་བལྟ་རུ་བཅུག ཤིང་མཁན་གྱིས་སྒོ་དེའི་ཚལ་བ་དུ་མ་མཐོང་། ཤིང་མཁན་གྱིས་བུ་ལ་བསམས། ཤིང་ངན་པ་འདི་\voc[གཅོག་]{ཆོག་}{to knock out}ཅིག སྒོ་འདི་ལ་བལྟོས་ཤིག་སྙམ་མོ། བུས་སྒོ་ལ་བལྟས་ནས་དེ་ཀ་བཅག་གོ།

དེའི་ཚེ་གྲོང་ཁྱེར་དུ་མི་ངན་པ་ཞིག་འོངས། མི་དེས་ཁྱིམ་བདག་གི་གོས་དང་ནོར་མང་པོ་\voc[རྐུ་]{བརྐུས་}{to steal}སོ། ཁྱིམ་བདག་གི་བུ་གཅིག་གིས་ནོར་དེ་\voc[འབོར་]{བོར་}{to lose}བ་མཐོང་ནས། བདག་གིས་ནོར་བོར་རོ། འདི་རྐུ་བ་ངན་པ་ཡིན་ནོ་ཞེས་\voc[སྨྲ་]{སྨྲས་}{to speak}སོ། རྐུ་བ་དེ་ཚལ་དུ་སོང་། བུ་དེས་ཤིང་མཁན་དང་ལྷན་ཅིག་རྐུ་བ་ཕྱིར་ཁྲིད་དེ་ཁྱིམ་གྱི་སྒོར་བཞག་གོ། བགྲེས་པས་རྐུ་བ་ལ་མ་སྤྱོས། བགྲེས་པས་དེ་སྐད་སྨྲས། ནོར་རྐུ་བ་ཉེས་པ་ཡིན། འོན་ཀྱང་སེམས་འཇམ་པར་གྱིས་ཤིག ཆོས་ཀྱི་\voc{དབང་}{power}གིས་ཉེས་པ་བཟང་པོར་འགྱུར་སྲིད་དོ་ཞེས་སོ།

རྐུ་བ་དེ་ནི་ལྟོགས་པ་ཡིན། ཁྱིམ་བདག་གི་ཨ་མས་རྐུ་བ་དེ་ལ་\voc{སྨན་}{medicine}དང་ཟས་དྲོན་མོ་བྱིན། རྐུ་བས་ཟས་བཟས་ནས་\voc{ཚིམ་པ་}{satisfied}བྱུང་། དེ་ནས་རྐུ་བས་སྨྲས། བདག་གི་\voc{ལས་}{action}ངན་པ་ཡིན། བདག་གིས་ནོར་བརྐུས་པ་ཉེས་པ་ཡིན། ད་ནས་བདག་གིས་ནོར་མི་རྐུའོ་ཞེས་སོ། བགྲེས་པས་རྐུ་བ་དེ་ལ་སྨྲས། དེ་ལྟ་བུའི་སྨྲས་པ་བཟང་པོ་ཡིན། ཁྱོད་ཀྱིས་སེམས་གསལ་པར་གྱིས་ཤིག་ཅེས་སོ།

དེ་ནས་གྲོང་ཁྱེར་དེར་\voc{ཇོ་བོ་}{lord}བཙུན་པ་ཞིག་\voc{གཤེགས་}{to come}སོ། ཇོ་བོ་དེ་ནི་\voc{བཙུན་པ་}{noble}དང་ཐུགས་རྗེ་ཅན་ཡིན། ཁྱིམ་བདག་དང་མ་སྨད་དང་བུ་ལྔ་ཀ་དང་ཕྲུ་གུ་དྲུག་ཀས་ཇོ་བོ་ལ་\voc{གསོལ་}{to request}བ། ཇོ་བོ། བདག་ཅག་ལ་ཆོས་སྟོན་ཅིག རྐུ་བ་དེ་ལ་ཡང་ཆོས་སྟོན་ཅིག་ཅེས་སོ། ཇོ་བོས་དེ་ཀུན་ལ་འཁོར་བ་དང་སྐྱེ་བ་དང་ལས་ཀྱི་དོན་བསྟན། དེའི་ཚིག་ཐོས་པ་ཀུན་གྱི་སེམས་ཚིམ་པར་གྱུར་ཏོ།

གྲོང་ཁྱེར་དེའི་ཕྱིར་ན་སྲིན་མོ་ངན་པ་ཞིག་ཡོད། སྲིན་མོ་དེའི་\voc{མདུན་སོ་}{incisor}དྲུག་རྣོན་པོ་ཡོད། སྲིན་མོ་དེས་\voc{བྱ་}{bird}དང་ཕུག་རོན་མང་པོ་བསད། ཁྱིམ་བདག་གི་བུས་བྱ་ཞིག་ཁྱིམ་དུ་ཁྲིད་དེ་སྨན་བྱིན། ཇོ་བོས་སྲིན་མོ་ལ་ཆོས་བསྟན། སྲིན་མོས་ཇོ་བོའི་ཚིག་ཐོས་ནས་སྨྲས། བདག་གིས་ལས་ངན་པ་བྱས། ད་ནས་བྱ་མི་གསོད་དོ་ཞེས་སོ། ཇོ་བོས་སྲིན་མོ་ལ་སྨན་དང་ཆོས་བྱིན་པས། སྲིན་མོའི་སེམས་གསལ་པར་གྱུར་ཏོ།

དེའི་ཚེ་རྒྱལ་པོ་དེའི་\voc{ལྷ་མོ་}{queen}ཡང་གནས་དེར་འོངས། ལྷ་མོ་དེ་ནི་བུ་མོ་དང་ལྷན་ཅིག་གཤེགས། ལྷ་མོས་ཇོ་བོ་ལ་གསོལ་བ་བྱས་ནས་སྨྲས། ཇོ་བོ། དེ་སྐད་སྨྲས་པའི་དོན་གསལ་པོ་ཡིན། སེམས་ཅན་ཀུན་ལ་ཆོས་སྟོན་པར་གསོལ། རྐུ་བ་དང་སྲིན་མོ་ཡང་སེམས་ཅན་ཡིན། དེ་ཀོ་བཟང་པོར་བྱེད་པ་ཆོས་ཀྱི་དབང་ཡིན་ནོ་ཞེས་སོ། ཇོ་བོས་ལྷ་མོའི་ཚིག་ཐོས་ནས་ཐུགས་དགའོ། དེ་ནས་ཁྱིམ་བདག་དང་བགྲེས་པ་དང་མ་སྨད་དང་ཤིང་མཁན་དང་རྐུ་བ་དང་སྲིན་མོ་ཀུན་གྱིས་ཆོས་ཐོས་ནས་སེམས་གསལ་བ་དང་ཚིམ་པ་བྱུང་ངོ་།
\end{readersection}

\begin{readersection}{མཉན་དུ་ཡོད་པ་}
\voc{མཉན་དུ་ཡོད་པ་}{Śrāvasti}ཞེས་བྱ་བའི་གྲོང་ཁྱེར་ཆེན་པོ་ན་\voc{སངས་རྒྱས་}{Buddha}འཁོར་མང་པོ་དང་བཅས་ཏེ་\voc{བཞུགས་}{to stay}སོ། གྲོང་ཁྱེར་དེའི་\voc{ཕོ་བྲང་}{palace}གི་སྒོ་ན་\voc{གཉེན་ཚན་}{relatives}མང་པོ་དང་\voc{བུ་སྨད་}{family}དང་\voc{ཚོགས་}{troop}མང་པོ་འདུག དེ་ཀུན་གྱིས་སངས་རྒྱས་ལ་ཆོས་ཉན་དུ་སོང་ངོ་། དེར་\voc{ཁོ་བོ་}{I}ཞེས་སྨྲ་བའི་ཁྱིམ་བདག་གཅིག་ཡོད། ཁྱིམ་བདག་དེའི་\voc{ཁྱིམ་}{house}ན་ཕ་མ་དང་བུ་ཚ་དང་ནོར་མང་པོ་ཡོད། དེའི་\voc{སོམ་ཚུགས་}{chin}རིང་པོ་དང་ཁ་སྤུ་རོག་པོ་ཡོད། དེས་\voc{འཁར་བ་}{staff}བཟུང་སྟེ་ཕོ་བྲང་གི་སྒོར་སོང་།

དེའི་ཚེ་\voc{གོས་སེར་ཅན་}{Viṣṇu}གྱི་གཏམ་ཐོས་པའི་བྲམ་ཟེ་ཞིག་དེར་འོངས། བྲམ་ཟེ་དེས་\voc{མཁའ་ལྡིང་}{garuḍa}ནམ་མཁའ་ལ་འཕུར་བ་མཐོང་ནས། མཁའ་ལྡིང་ནི་བྱ་ཆེན་པོ་ཡིན་ནོ་སྙམ། མཁའ་ལྡིང་གི་གཟུགས་ཆེ་བ་དང་སྐད་དྲག་པ་མཐོང་སྟེ། སྤེའུ་ཕྲུག་དང་ཕུག་རོན་ཀུན་ཏུ་འཇིགས། བྲམ་ཟེས་སྨྲས། གོས་སེར་ཅན་གྱི་དབང་གིས་མཁའ་ལྡིང་འཕུར་རོ་ཞེས་སོ། དེ་སྐད་ཐོས་ནས་ཁྱིམ་བདག་གིས། བདག་ཅག་གི་སངས་རྒྱས་ནི་འཁོར་བའི་དག་བོ་འཇོམས་པར་མཁྱེན་ནོ་སྙམ་མོ།

ཕོ་བྲང་གི་ཕྱིར་ན་\voc{དག་བོ་}{enemy}མང་པོ་བྱུང་། དག་བོ་དེ་ཀུན་གྱིས་\voc{རྒོད་མ་}{mare}དང་\voc{བཞོན་པ་}{riding-beast}མང་པོ་བཟུང་ནས་ཕོ་བྲང་གི་སྒོར་སོང་། རྒྱལ་པོ་དེ་གཅིག་པུ་མ་ཡིན། འཁོར་དང་བུ་སྨད་དང་གཉེན་ཚན་མང་པོ་ཡོད། འོན་ཀྱང་\voc{གཅིག་པུ་}{sole}བུ་ཐ་ཆུངས་དེ་སྒོ་ན་འདུག དག་བོ་གཅིག་གིས་རྟ་ལ་འཁར་བ་\voc[འཆའ་]{བཅས་}{to place}ཏེ། རྟ་དེ་བཞོན་པར་བྱས་སོ། དག་བོ་གཞན་གྱིས་རྒོད་མ་བཟུང་ནས་ཕྲིན་བྱེད་པར་བསམས།

རྒྱལ་པོས་དེ་མཐོང་ནས་ཞལ་ཆེ་གཅོད་པར་འདོད། རྒྱལ་པོས་འཁོར་ལ་སྨྲས། དག་བོ་འདི་ཀུན་གྱིས་ནོར་བརྐུས། དེས་ན་ཁོང་ལ་\voc{ཞལ་ཆེ་}{judgement}\voc{གཅོད་}{to decide}པར་བྱའོ་ཞེས་སོ། ཁྱིམ་བདག་དེས་སྨྲས། དག་བོ་ཀུན་མེད་པར་བྱས་ན་སྐྱིད་སྡུག་རྒྱས་པར་མི་འགྱུར་རམ་ཞེས་སོ། རྒྱལ་པོས་དེ་སྐད་ཐོས་ནས་སྨྲས། ཁོ་བོས་དག་བོ་འདི་ཀུན་\voc[འགུམས་]{བཀུམ་}{to kill}པར་མི་བྱ། སངས་རྒྱས་ལ་ཞལ་ཆེ་གཅད་པར་གསོལ་ལོ་ཞེས་སོ།

དེ་ནས་རྒྱལ་པོ་དང་ཁྱིམ་བདག་དང་དག་བོ་ཀུན་སངས་རྒྱས་ཀྱི་གནས་སུ་སོང་། སངས་རྒྱས་བཞུགས་པ་ལ་ཀུན་གྱིས་ཕྱག་བྱས། རྒྱལ་པོས་དག་བོའི་ལས་ངན་པ་བསྟན། དག་བོ་ཀུན་གྱིས་བདག་ཅག་གི་ལས་ངན་པ་ཡིན་ནོ་ཞེས་སྨྲས། སངས་རྒྱས་ཀྱིས་དག་བོ་ལ་བལྟས་ཏེ། མི་གསོད་པ་དང་ནོར་མི་རྐུ་བ་དང་ཚིག་ངན་པ་མི་སྨྲ་བ་ནི་\voc{བསྟན་པ་}{teaching}བཟང་པོ་ཡིན་ནོ་ཞེས་ཆོས་བསྟན། དག་བོ་ཀུན་གྱིས་དེ་ཐོས་ཏེ་སེམས་འཇམ་པར་གྱུར།

དེར་ཁྱིམ་བདག་དེས། \voc{ད་ལྟར་}{now}ཁོ་བོ་ལ་སེམས་གསལ་པོ་སྐྱེས་སོ། འཁོར་བའི་\voc{སྐྱིད་སྡུག་}{joy and sorrow}ཀུན་ཏུ་འཁོར་ཏེ་འགྲོ། སེམས་ཅན་གྱི་ལས་ཀྱིས་བསྟན་པ་\voc{དར་}{to spread}རམ་འཛད་པར་འགྱུར། ཆོས་དར་ན་སེམས་ཅན་ཀུན་གྱི་ཡིད་རངས་པར་འགྱུར་རོ་སྙམ་མོ། དེ་ནས་ཁྱིམ་བདག་གིས་སངས་རྒྱས་ལ་གསོལ་བ་བྱས་ཏེ། བདག་ཅག་གི་ཁྱིམ་དུ་ཕྲིན་སྐུར་བར་འདོད། གཉེན་ཚན་དང་བུ་སྨད་ཀུན་ལ་བསྟན་པ་བསྟན་པར་གསོལ་ཞེས་སོ།

སངས་རྒྱས་ཀྱིས་ཀུན་དགའ་བོ་ལ་བཀའ་བསྩལ་ཏེ། ཁྱིམ་དེར་\voc{ཕྲིན་}{message}\voc{བྱེད་}{to say}དུ་སོང་ཞིག གཉེན་ཚན་ཀུན་ལ་ཆོས་བསྟན་པར་བྱོས་ཤིག་ཅེས་སོ། ཀུན་དགའ་བོས་ཕྲིན་\voc[བྱེད་]{བྱོས་}{to say}ཏེ་ཁྱིམ་དེར་སོང་། ཁྱིམ་ན་བུ་ཚ་དང་མ་སྨད་དང་རྒན་མོ་དུ་མ་འདུག ཀུན་དགའ་བོས་ཁོ་བོ་སངས་རྒྱས་ཀྱི་ཕྲིན་བྱེད་དུ་འོངས་སོ། བསྟན་པ་ཟབ་མོ་ཉོན་ཅིག་ཅེས་སྨྲས། དེ་ཀུན་གྱིས་ཕྲིན་ཐོས་ཏེ་ཡི་རངས་སོ།

དག་བོ་དེ་ཀུན་ན་གཞོན་གཅིག་པུ་ཞིག་ཡོད། དེའི་\voc{ཡི་}{mind}ངན་པ་རྒྱས་ནས། སླར་ཡང་དག་བོའི་ཚོགས་དང་ལྷན་ཅིག་ནོར་རྐུ་བར་བསམས། དེས་རྒོད་མ་བཟུང་ནས་ཕོ་བྲང་གི་སྒོར་འོངས། སངས་རྒྱས་ཀྱིས་དེ་མཁྱེན་ཏེ་འཁོར་ལ་སྨྲས། གཞོན་པ་འདི་གཞན་དུ་བསམས། འོན་ཀྱང་གསོད་པར་མ་བྱ། ཆོས་ཀྱིས་དག་བོ་\voc[འཇོམས་]{བཅོམ་}{to conquer}པར་བྱའོ་ཞེས་སོ། སངས་རྒྱས་ཀྱི་བསྟན་པས་དེའི་ཡི་ཞི་བར་གྱུར་ཏེ། དེས་ནོར་མ་བརྐུས་སོ།

གྲོང་ཁྱེར་དེར་ཆོས་ཀྱི་སྐད་ཀུན་ཏུ་དར་རོ། ཕོ་བྲང་ནས་ཁྱིམ་གྱི་བར་དུ་མི་ཀུན་དགའ། ཁྱིམ་བདག་གི་བུ་སྨད་ཀྱིས་ཟས་དང་གོས་སེམས་ཅན་ལ་བྱིན། སྐྱིད་སྡུག་མི་རྟག་པར་ཤེས་ཏེ། ནོར་\voc{ཙམ་}{mere}ལ་མ་ཆགས། མཁའ་ལྡིང་ནམ་མཁའ་ལ་འཕུར་བ་དང་རྒོད་མ་ལམ་དུ་འགྲོ་བ་དང་བྱ་རྩེ་བ་ཀུན་གྱིས་མཐོང་། བུ་ཚ་དེ་ཀུན་གྱིས་ཚལ་ན་\voc{རྩེ་}{to play}ཞིང་ཆོས་ཀྱི་ཚིག་ཐོས་སོ།

དེའི་ཚེ་ཁྱིམ་བདག་གིས་སངས་རྒྱས་ལ་སྨྲས། \voc{འུ་བུ་ཅག་}{we}གིས་ད་ལྟར་བསྟན་པའི་དོན་\voc{ཤེས་}{to know}སོ། བདག་ཅག་གིས་ནོར་གཏེར་དུ་མི་\voc{འཛིན་}{to take}། ནོར་ནི་སེམས་ཅན་ལ་སྦྱིན་པར་བྱ། དག་བོ་ཡང་ཆོས་ཀྱིས་འཇོམས་པར་བྱ། སེམས་ཅན་གསོད་པར་མི་བྱ། དེ་ལྟར་བསམས་ཏེ་ཁོ་བོའི་ཡི་རངས་སོ་ཞེས་སོ། སངས་རྒྱས་ཀྱིས་དེ་ལ་བལྟས་ནས། དེ་ལྟ་བུའི་\voc{རང་}{to rejoice}བ་ནི་བསམ་པ་བཟང་པོ་ཡིན་ནོ་ཞེས་སོ། དེ་ཐོས་ནས་འཁོར་དང་མི་ཀུན་གྱི་\voc[སེམས་]{བསམས་}{to think}པ་རྒྱས་ཤིང་སྐྱིད་སྡུག་གི་དོན་ཤེས་ཏོ།
\end{readersection}

\begin{readersection}{བྱ་ཚོགས་ཀྱི་ནགས་མ་}
\voc{མཉན་དུ་ཡོད་པ་}{Śrāvasti}གྲོང་ཁྱེར་དུ་\voc{གཽ་ཏ་མ་}{Gautama}སངས་རྒྱས་འཁོར་དགེ་སློང་མང་པོ་དང་བཅས་ཏེ་བཞུགས་སོ། དེའི་ཚེ་གྲོང་ཁྱེར་དེའི་ཕྱིར་\voc{བྱ་ཚོགས་}{Bya-chogs}\voc{བྱ་བ་}{so-called}འི་\voc{ནགས་མ་}{forest}ཞིག་ཡོད། ནགས་མ་དེའི་\voc{ནང་གི་}{among}\voc{ཤ་ཀོ་ཏ་ཀ་}{śākhoṭaka}ཞེས་བྱ་བའི་\voc{ཤིང་}{tree}ཆེན་པོ་ཞིག་འདུག ཤིང་དེའི་རྩེ་ན་\voc{ཁྲ་}{falcon}གཅིག་དང་\voc{བྱ་རོག་}{crow}གཅིག་འདུག ཤིང་དེའི་དྲུང་ན་\voc{བ་ར་བོང་}{cattle}དང་\voc{ཕྱུགས་}{cattle}མང་པོ་འདུག ནགས་མ་དེར་ཟས་དང་ཆུ་\voc{མོད་}{abundant}པས། སེམས་ཅན་\voc{ཐམས་ཅད་}{all}\voc{བདེ་སྐྱིད་}{to be happy}དོ།

གྲོང་ཁྱེར་དེར་\voc{ཟས་གཙང་མ་}{Zas-gcaṅ-ma}\voc{བཙུན་མོ་}{queen}དང་\voc{སྒྱུ་མ་ལྷ་མཛེས་}{Sgyu-ma Lha-mdzes}བུ་མོ་ཡོད་དོ། བཙུན་མོ་དེའི་\voc{ཡུམ་}{mother}ནི་རྒན་མོ་སྙིང་བརྩེ་བ་ཞིག་ཡིན། བཙུན་མོ་ལ་\voc{སྙིང་རྗེ་}{compassion}ཆེན་པོ་ཡོད། སྒྱུ་མ་ལྷ་མཛེས་ནི་བུ་མོ་\voc{ཆུང་ངུ་}{youngest}ཡིན་ཡང་། སེམས་\voc{སྙིང་བརྩེ་བ་}{kind}དང་ལྡན་ནོ། དེ་གཉིས་ཀྱིས་གྲོང་ཁྱེར་ནས་ཟས་དང་ཆུ་མང་པོ་\voc[སྐྱེལ་]{བསྐྱལ་}{to bring}ཏེ་བྱ་ཚོགས་ཀྱི་ནགས་མར་སོང་། བཙུན་མོས་ཟས་དེ་བླངས་ནས་བ་ར་བོང་ལ་\voc[སྙོད་]{བསྙོད་}{to feed}། སྒྱུ་མ་ལྷ་མཛེས་ཀྱིས་ཆུ་བླངས་ནས་ཕྱུགས་དང་བྱ་རོག་ལ་བསྙོད་དོ།

ནགས་མ་དེའི་ཕྱིར་ན་མི་ངན་པ་\voc{དྲུག་པོ་}{six}ཞིག་ཡོད། མི་དྲུག་པོ་དེ་ནི་ཡུལ་དེའི་\voc{དགྲ་}{enemy}ཡིན། དགྲ་དྲུག་པོས་\voc{རྡོ་}{stone}དང་འཁར་བ་བཟུང་སྟེ། བྱ་རོག་དང་ཁྲ་ལ་\voc[འཕེན་]{འཕངས་}{to throw}སོ། ཁྲ་དེ་ནམ་མཁའ་ལ་འཕུར་ཏེ་ཕྱིར་སོང་། བྱ་རོག་དེ་ཤིང་ལས་བྱི་སྟེ་ཆུར་ཚུད་དོ། སྒྱུ་མ་ལྷ་མཛེས་ཀྱིས་བྱ་རོག་དེ་\voc[འཁྱེར་]{ཁྱེར་}{to carry}ནས་གནས་བཟང་པོར་བཞག བཙུན་མོས་དགྲ་དེ་ཀུན་ལ་མ་སྤྱོས། དེས་སྙིང་རྗེས་སྨྲས་པ། སེམས་ཅན་ཆུང་ངུ་ལ་རྡོ་མ་འཕང་ཞིག སེམས་ཅན་ཐམས་ཅད་ཀྱིས་བདེ་སྐྱིད་འདོད་དོ་ཞེས་སོ།

དགྲ་དྲུག་པོ་དེས་བཙུན་མོའི་ཚིག་ཐོས་ཀྱང་སེམས་མ་འཇམ། དེ་ཀུན་གྱིས་བ་ར་བོང་གཅིག་དང་ཕྱུགས་གཉིས་\voc[ལེན་]{བླངས་}{to take}ནས་ནགས་མའི་ནང་དུ་སོང་། བཙུན་མོའི་ཡུམ་གྱིས་དེ་མཐོང་ནས། དགྲ་དྲུག་པོ་ནི་སེམས་ངན་པ་དང་བཅས་པ་ཡིན། འོན་ཀྱང་ཆོས་ཐོས་ན་སེམས་བཟང་པོར་འགྱུར་སྲིད་དོ་སྙམ་མོ། དེ་ནས་ཡུམ་དེས་དགེ་སློང་ཞིག་ལ་བརྡ་བསྟན། དགེ་སློང་དེས་སངས་རྒྱས་ལ་ཕྲིན་བྱས་སོ།

དེ་ནས་སངས་རྒྱས་ནགས་མ་དེར་གཤེགས་ཏེ། ཤ་ཀོ་ཏ་ཀའི་ཤིང་དྲུང་ན་བཞུགས་སོ། \voc{དགེ་སློང་}{monk}མང་པོ་དང་ཁྱིམ་བདག་དང་བཙུན་མོ་དང་སྒྱུ་མ་ལྷ་མཛེས་ཐམས་ཅད་དེར་འདུག སངས་རྒྱས་ཀྱིས་དགྲ་དྲུག་པོ་ལ་བལྟས་ཏེ། གསོད་པ་དང་རྐུ་བ་དང་རྡོ་འཕེན་པ་ནི་ལས་ངན་པ་ཡིན། སྙིང་རྗེ་དང་སྙིང་བརྩེ་བ་ནི་ཆོས་བཟང་པོ་ཡིན། བྱང་ཆུབ་སེམས་དཔའས་སེམས་ཅན་ཐམས་ཅད་ལ་ཟས་བསྙོད་ཅིང་ཆུ་སྐྱེལ་ལོ་ཞེས་བསྟན་ཏོ།

དེའི་ཚེ་དགྲ་དྲུག་པོའི་ནང་གི་ཆུང་ངུ་ཞིག་གིས་དེ་སྐད་ཐོས་ནས་ངུས་སོ། དེས་རྡོ་ས་ལ་བཞག་སྟེ། བདག་གིས་སེམས་ཅན་ལ་རྡོ་འཕངས། བདག་གི་ལས་ངན་པ་ཡིན། ད་ནས་སེམས་ཅན་ལ་མི་འཕེན་ནོ་ཞེས་སྨྲས། གཞན་དག་གིས་ཀྱང་བ་ར་བོང་དང་ཕྱུགས་ཕྱིར་ཁྱེར་ནས་བཙུན་མོ་ལ་བྱིན་ནོ། བཙུན་མོས་དགྲ་དེ་ཀུན་ལ་ཟས་གཙང་མ་བསྙོད། སྒྱུ་མ་ལྷ་མཛེས་ཀྱིས་ཆུ་བསྐྱལ་ཏེ་ཁོང་ཀུན་ལ་བྱིན་ནོ།

ཁྲ་དེ་སླར་འོངས་ཏེ་ཤིང་གི་རྩེ་ན་འདུག བྱ་རོག་ཀྱང་ཆུ་ནས་ཐར་ཏེ་ཤིང་དྲུང་ན་འདུག བ་ར་བོང་དང་ཕྱུགས་ཐམས་ཅད་ཀྱང་བདེ་སྐྱིད་དོ། སངས་རྒྱས་ཀྱིས་བྱ་རོག་དང་ཁྲ་དང་བ་ར་བོང་ལ་བལྟས་ཏེ། སེམས་ཅན་ཆུང་ངུ་ཡང་ཚེ་སྲོག་ལ་ཆགས། དེས་ན་སྙིང་རྗེ་ཆེན་པོ་སྐྱེད་ཅིག་ཅེས་སོ། དགེ་སློང་ཐམས་ཅད་ཀྱིས་ཆོས་དེ་ཐོས་ཏེ་ཡི་རངས་སོ།

དེ་ནས་སྒྱུ་མ་ལྷ་མཛེས་ཀྱིས་སངས་རྒྱས་ལ་གསོལ་བ། བདག་ནི་བུ་མོ་ཆུང་ངུ་ཡིན་ཡང་། སེམས་ཅན་ཐམས་ཅད་ལ་ཟས་སྐྱེལ་བར་འདོད། བདག་གི་\voc{ཕ་}{father}དང་ཡུམ་དང་བཙུན་མོས་ཀྱང་སྙིང་རྗེ་བསྟན། བདག་ཅག་གིས་ནགས་མ་དེར་ཕྱུགས་དང་བྱ་ཐམས་ཅད་བསྙོད་པར་བྱའོ་ཞེས་སོ། སངས་རྒྱས་ཀྱིས་དེ་ལ་སྨྲས་པ། བུ་མོ་ཆུང་ངུ་ཡང་སྙིང་བརྩེ་བ་ཡོད་ན། \voc{བྱང་ཆུབ་སེམས་དཔའ་}{bodhisattva}ཡི་ལས་བྱེད་དོ་ཞེས་སོ།

དེའི་ཚེ་དགྲ་དྲུག་པོ་ལས་གཅིག་གི་སེམས་སླར་ངན་པར་གྱུར། དེས་ནོར་དང་ཟས་\voc[ལེན་]{བླངས་}{to take}ནས་\voc[འབྲོས་]{བྲོས་}{to flee}སོ། དགེ་སློང་ཞིག་གིས་དེ་མཐོང་ནས་ཕྱིར་ཁྲིད་དོ། སངས་རྒྱས་ཀྱིས་དེ་ལ་མ་སྤྱོས་ཏེ། ཁྱོད་ཀྱིས་ཟས་བླངས་ནས་བྲོས་པ་ནི་ཉེས་པ་ཡིན། འོན་ཀྱང་སེམས་བསྒྱུར་ན་ཆོས་ཀྱི་སྒོ་ཡོད་དོ་ཞེས་སོ། དགྲ་དེས་ཟས་ཕྱིར་བསྐྱལ་ཏེ། བདག་གིས་ཕྱིན་ཆད་མི་རྐུ། སེམས་ཅན་ལ་ཟས་བསྙོད་པར་བྱའོ་ཞེས་སོ།

དེ་ནས་མཉན་དུ་ཡོད་པའི་མི་ཐམས་ཅད་བྱ་ཚོགས་ཀྱི་ནགས་མར་འོངས་ཏེ། ཟས་དང་ཆུ་སེམས་ཅན་ལ་བསྐྱལ། དགྲ་དྲུག་པོ་ཡང་སེམས་འཇམ་པར་གྱུར་ཏེ། ཕྱུགས་བསྙོད་ཅིང་རྡོ་མ་འཕངས། ཁྲ་དང་བྱ་རོག་དང་བ་ར་བོང་ཐམས་ཅད་བདེ་སྐྱིད་པར་འདུག སངས་རྒྱས་ཀྱི་བསྟན་པ་ནགས་མ་དེར་རབ་ཏུ་དར། མི་དང་ཕྱུགས་དང་བྱ་ཐམས་ཅད་ལ་སྙིང་རྗེ་བྱུང་ངོ་།
\end{readersection}

\begin{readersection}{རིག་པའི་སྟེང་དུ་མཆོངས་པ་}
མཉན་དུ་ཡོད་པའི་གྲོང་ཁྱེར་ན་སངས་རྒྱས་འཁོར་དགེ་སློང་མང་པོ་དང་བཅས་ཏེ་བཞུགས་སོ། དེའི་ཚེ་གྲོང་ཁྱེར་གྱི་དབུས་ན་ཕོ་བྲང་ཆེན་པོ་ཞིག་ཡོད། ཕོ་བྲང་དེའི་སྒོ་ལ་འཁོར་མང་པོས་\voc[སྐོར་]{བསྐོར་}{to surround}ཏོ། དེར་ཁྱིམ་བདག་དབུལ་བ་ཞིག་འོངས། ཁྱིམ་བདག་དེ་ལ་\voc{བགོ་བ་}{clothes}ངན་པ་ཙམ་ཡོད། ཟས་མེད་པའི་ཕྱིར་ལུས་ལ་\voc{ངལ་}{weariness}ཆེན་པོ་བྱུང་། དེའི་བུ་ཚ་དྲུག་ནི་འགྲོ་བ་དབུལ་ཕོངས་ཡིན། ཁྱིམ་བདག་དེས་ཕོ་བྲང་གི་སྒོ་ནས་སངས་རྒྱས་ལ་བལྟས་ཏེ། སངས་རྒྱས་ནི་འགྲོ་བ་ཀུན་ལ་བྱམས་པ་དང་ལྡན་པ་ཡིན་ནོ་སྙམ་མོ།

དེ་ནས་ཁྱིམ་བདག་དེས་སངས་རྒྱས་ཀྱི་དྲུང་དུ་སོང་ནས་ཕྱག་བྱས་ཏེ། སངས་རྒྱས་ལ་གསོལ་བ། བདག་དབུལ་ཕོངས་ཡིན། བདག་ལ་ཟས་མེད། བགོ་བ་མེད། བུ་ཚ་དྲུག་ཀྱང་ལྟོགས་པ་ཡིན། \voc{ཅི་}{what}བྱ། \voc{ཅིའི་སླད་དུ་}{why}དབུལ་ཕོངས་འདི་འབྱུང་ཞེས་གསོལ་ཏོ། སངས་རྒྱས་ཀྱིས་ཁྱིམ་བདག་དེ་ལ་ཆོས་བསྟན་ནས། འགྲོ་བ་ན་སྐྱེ་བ་དང་འཇིག་པ་དང་དབུལ་ཕོངས་དང་བདེ་བ་ཡོད། དེ་ཀུན་ལས་ཀྱི་ཕྱིར་འབྱུང་ངོ་ཞེས་གསུངས། དགེ་སློང་ཀུན་གྱིས་སངས་རྒྱས་ལ་\voc[མངོན་པར་སྟོད་]{མངོན་པར་བསྟོད་}{to praise}དོ།

ཕོ་བྲང་གི་དྲུང་ན་\voc{རིག་པ་}{wall}རིང་པོ་ཞིག་ཡོད། རིག་པ་དེའི་\voc{སྟེང་དུ་}{on}བྱ་རོག་དང་ཁྲ་གཉིས་འདུག ཁྲ་དེས་བྱ་རོག་ལ་བལྟས་ནས་སྐད་དྲག་པོ་བྱས། བྱ་རོག་དེ་འཇིགས་ཏེ་རིག་པའི་སྟེང་ནས་མཆོངས། དེ་ལྟར་\voc{མཆོངས་}{to jump}པའི་ཕྱིར་རྐང་པ་ཅུང་ཟད་ཆག་གོ། ཁྱིམ་བདག་གི་བུ་ཆུང་ངུས་བྱ་རོག་དེ་མཐོང་ནས་བཟུང་སྟེ་སངས་རྒྱས་ཀྱི་དྲུང་དུ་ཁྱེར། བུ་དེས་འདི་སྐད་ཅེས་གསོལ། བྱ་རོག་འདི་རིག་པའི་སྟེང་ནས་མཆོངས་ཏེ་རྐང་པ་ཆག འདི་ལ་\voc{སྙུན་}{illness}ཡོད་དམ་ཞེས་སོ།

སངས་རྒྱས་ཀྱིས་བྱ་རོག་ལ་བལྟས་ཏེ། འདི་ནི་སྙུན་ཆེན་པོ་མིན། རྐང་པ་ཆག་པའི་ཕྱིར་མྱ་ངན་ཅུང་ཟད་ཚོར་རོ་ཞེས་གསུངས། དེ་ནས་སངས་རྒྱས་ཀྱིས་དགེ་སློང་ཞིག་ལ། བྱ་རོག་འདི་ལ་\voc[གསོ་]{བསོ་}{to cure}བར་བྱོས་ཤིག ངལ་ཆེན་པོ་ཡོད་པས་\voc[གསོ་]{ངལ་གསོ་}{to cure}བར་བྱོས་ཤིག་ཅེས་གསུངས། དགེ་སློང་དེས་བྱ་རོག་ལ་ཟས་དང་ཆུ་བྱིན། བྱ་རོག་དེས་\voc{ཟས་}{food}ཟོས་ནས་ངལ་གསོས་སོ།

དེའི་\voc{བར་དུ་}{during}ཕོ་བྲང་ནས་མི་ལྔ་དང་བུ་མོ་གཅིག་འོངས། མི་དེ་ཀུན་གྱིས་\voc{འདི་སྐད་}{these words}གསོལ། དབུལ་ཕོངས་འདི་ལ་ཟས་དང་བགོ་བ་མང་པོ་སྦྱིན་པར་བགྱི། འོན་ཀྱང་དབུལ་ཕོངས་ཀྱི་\voc{ཕྱིར་}{on account of}ཁོང་གི་སེམས་ལ་མྱ་ངན་ཆེན་པོ་ཡོད། ཅི་བགྱི་ཞེས་གསོལ། སངས་རྒྱས་ཀྱིས། ཟས་སྦྱིན་པ་དང་བགོ་བ་སྦྱིན་པ་ནི་བཟང་། འོན་ཀྱང་ཆོས་སྟོན་པ་ནི་དེ་བས་ཀྱང་བཟང་ངོ་ཞེས་གསུངས།

དེ་ནས་སངས་རྒྱས་ཀྱིས་ཁྱིམ་བདག་དེ་ལ་སྐྱེ་བ་དང་འགྲོ་བ་དང་ལས་ཀྱི་དོན་བསྟན། ཁྱིམ་བདག་དེས་ཆོས་ཐོས་ནས་མྱ་ངན་ཅུང་ཟད་ཟད། དེས་གསོལ་བ། བདག་ནི་དབུལ་ཕོངས་\voc{འབའ་ཞིག་}{alone}ཏུ་མ་ཟད། བདག་གི་སེམས་ལ་མྱ་ངན་ཡང་ཡོད། ད་ནས་བདག་གིས་ཆོས་ཉན་པར་བྱའོ་ཞེས་སོ། སངས་རྒྱས་ཀྱིས། ཆོས་ཉན་པ་ནི་མྱ་ངན་གསོ་བའི་སྨན་ཡིན་ནོ་ཞེས་གསུངས།

དེའི་ཚེ་\voc{མི་ལ་ཤེས་རབ་རྒྱལ་མཚན་}{Mi-la father}ཞེས་བྱ་བའི་ཁྱིམ་བདག་གཞན་ཞིག་དེར་འོངས། དེའི་རིགས་ནི་བཟང་པོ་ཡིན་ཡང་། དེའི་ལུས་ལ་\voc{སྙུན་}{illness}ཆེན་པོ་ཡོད། དེས་སངས་རྒྱས་ལ་གསོལ་བ། བདག་གི་སྙུན་གྱི་ཕྱིར་ངལ་ཆེ། བདག་གི་བུ་དང་བུ་མོ་ལ་ཟས་དང་བགོ་བ་ཡོད་ཀྱང་། བདག་གི་ཚེ་འཕོ་བར་འགྱུར་རམ། བདག་གི་ལུས་དང་\voc{རུས་པ་}{bone}ཅིར་འགྱུར་ཞེས་གསོལ། སངས་རྒྱས་ཀྱིས། སྐྱེས་པ་ཀུན་ཚེ་འཕོ་བར་འགྱུར། ལུས་དང་རུས་པ་འཇིག་པར་འགྱུར། ཆོས་ཀྱི་དོན་ཤེས་པས་མྱ་ངན་ཉུང་ངུར་འགྱུར་རོ་ཞེས་གསུངས།

དེ་ནས་སངས་རྒྱས་ཀྱིས་རྫུ་འཕྲུལ་ཅུང་ཟད་བསྟན། \voc{རྫུ་འཕྲུལ་}{magic}དེས་མི་ཀུན་གྱི་མྱ་ངན་མ་བསལ་ཡང་། ཆོས་ཀྱི་སྒོ་བསྟན་ཏོ། དགེ་སློང་ཞིག་གིས་གསོལ་བ། རྫུ་འཕྲུལ་འདི་ཅིའི་སླད་དུ་བསྟན་ཞེས་སོ། སངས་རྒྱས་ཀྱིས། རྫུ་འཕྲུལ་འབའ་ཞིག་གིས་འགྲོ་བ་མི་སྐྱོབ། ཆོས་ཀྱི་དོན་ཚོར་བ་ནི་གསོ་བ་ཡིན། དེ་སྟོན་པའི་ཕྱིར་བསྟན་ཏོ་ཞེས་གསུངས། དེ་ཐོས་ནས་མི་ཀུན་གྱིས་ཆོས་ཀྱི་དོན་\voc{ཚོར་}{to perceive}རོ།

དེ་ནས་ཕོ་བྲང་གི་ནང་ནས་ཟས་དང་\voc{བཟའ་བ་}{food}མང་པོ་ཕྱུང་། རྒྱལ་པོས་དབུལ་ཕོངས་དང་སྙུན་ཅན་ཀུན་ལ་ཟས་དང་བཟའ་བ་བྱིན། ཁྱིམ་བདག་དབུལ་བ་དེ་ལ་ཟས་\voc{ཟིན་}{to be seized}ཏེ་ལུས་ཀྱི་ངལ་ཉུང་ངུར་གྱུར། མི་ལ་ཤེས་རབ་རྒྱལ་མཚན་ལ་ཡང་ཟས་བྱིན། དེས་ཟས་ཟོས་ཀྱང་། སྙུན་གྱི་ཕྱིར་མྱ་ངན་ཚོར་རོ། དེའི་བུ་ཚ་གསུམ་པོས་ཕ་ལ་བལྟས་ནས་ངུས་སོ།

སངས་རྒྱས་ཀྱིས་བུ་ཚ་\voc{གསུམ་པོ་}{three}ལ་གསུངས་པ། ཁྱེད་ཅིའི་སླད་དུ་ངུ་། \voc{སུ་}{who}ཡང་སྐྱེས་ནས་མི་འཆི་བ་མེད། སུ་ཡང་ལུས་དང་རུས་པ་ལ་རྟག་ཏུ་མི་གནས། ཁྱེད་ཀྱི་ཕས་ཆོས་ཐོས་ན་མྱ་ངན་ཉུང་ངུར་འགྱུར་རོ་ཞེས་གསུངས། བུ་ཚ་གསུམ་པོས་སངས་རྒྱས་ཀྱི་\voc[གསུང་]{གསུངས་}{to speak}པ་ཐོས་ནས། ཕའི་སྙུན་གྱི་བར་དུ་ཆོས་ཉན་པར་བྱ། མྱ་ངན་འབའ་ཞིག་ཏུ་མི་གནས་སོ་སྙམ་མོ།

དེར་\voc{ཤིང་རྟ་ཆེན་པོ་}{Śiṅ-rta-čhen-po}ཞེས་བྱ་བའི་རྒྱལ་བུ་ཡང་འོངས། དེས་སངས་རྒྱས་ལ་གསོལ་བ། བདག་གིས་སྔོན་གྱི་\voc{ལེའུ་}{chapter}གཅིག་བཀླགས། ལེའུ་དེར་རྒྱལ་པོས་འགྲོ་བ་དབུལ་ཕོངས་ལ་ཟས་སྦྱིན་པ་བསྟན་ཏོ། སྦྱིན་པ་དེ་ཅིའི་སླད་དུ་བྱས་ཞེས་གསོལ། སངས་རྒྱས་ཀྱིས་གསུངས་པ། འགྲོ་བ་ཀུན་ཟས་འདོད། འགྲོ་བ་ཀུན་མྱ་ངན་མི་འདོད། དེའི་ཕྱིར་སྦྱིན་པ་བྱའོ་ཞེས་སོ། ཤིང་རྟ་ཆེན་པོས་སངས་རྒྱས་མངོན་པར་བསྟོད་དོ།

དེའི་ཚེ་ཕོ་བྲང་གི་རིག་པའི་སྟེང་དུ་སླར་བྱ་རོག་དེ་འདུག ཁྲ་དེ་ཡང་དེར་འོངས། འོན་ཀྱང་ཁྲ་དེས་བྱ་རོག་ལ་མ་བསད། མ་འཇིགས་པར་གཅིག་གི་དྲུང་ན་གཅིག་འདུག མི་ཀུན་གྱིས་དེ་མཐོང་ནས་སྙིང་རྗེའི་དོན་ཚོར་རོ། སངས་རྒྱས་ཀྱིས་འདི་སྐད་གསུངས། གང་གིས་གཞན་ལ་ཉེས་པ་མི་བྱེད་པ་དེ་ནི་ཆོས་མཐོང་བ་ཡིན། གང་གིས་མྱ་ངན་ཚོར་བ་དེས་གཞན་གྱི་མྱ་ངན་ཡང་ཤེས་སོ། དེ་ཐོས་ནས་འཁོར་ཀུན་གྱི་སེམས་དགའོ།
\end{readersection}

\begin{readersection}{བོད་ཀྱི་ལོ་རྒྱུས་}
\voc{སྔོན་}{earlier}\voc{དུས་}{time}\voc{བོད་}{Tibet}ཀྱི་ཡུལ་ན་\voc{སྒྲ་ཆེན་པོ་}{Sgra-chen-po}ཞེས་བྱ་བའི་\voc{རྗེ་}{lord}ཞིག་ཡོད་དོ། རྗེ་དེའི་ཕོ་བྲང་\voc{ནང་དུ་}{in}བུ་སྨད་དང་གཉེན་ཚན་མང་པོ་བཞུགས། ཕོ་བྲང་གི་\voc{དྲུང་ན་}{under}ཤ་ཀོ་ཏ་ཀའི་ཤིང་ཆེན་པོ་ཞིག་སྐྱེས། ཤིང་དེ་\voc{སྤུ་}{hair}ལྟ་བུའི་འདབ་མ་མང་པོ་དང་\voc{སེ་འུ་}{pomegranate}དང་ལྡན་ནོ། སེ་འུ་དེའི་\voc{ཁ་དོག་}{colour}ནི་ངུར་སྨྲིག་དང་འདྲ། དེའི་\voc{བྲོ་བ་}{taste}ནི་\voc{ཅུང་ཟད་}{a little}བཟང་པོ་ཡིན་ཡང་། རྗེ་དེས་རང་\voc{ཉིད་}{very}ཀྱི་དོན་དུ་མ་བཟས། སེམས་ཅན་ལྟོགས་པ་ཐམས་ཅད་ལ་\voc[གཏོང་]{བཏང་}{to let go}ངོ་།

རྗེ་དེའི་སྲས་\voc{རབ་}{eldest}ནི་ཐུགས་བྱམས་པ་དང་\voc{སྤྲོ་བ་}{joy}དང་ལྡན་པ་ཡིན། \voc{གཅིག་པ་}{only}བུ་ཐ་ཆུངས་དེ་ནི་ཆུང་ངུའི་དུས་ནས་མགོ་ལ་\voc{སྤུ་}{hair}རིང་པོ་དང་ཁ་སྤུ་རོག་པོ་ཡོད། དེའི་\voc{མཇུག་མ་}{tail}མེད་ཀྱང་། བུ་ཆུང་ངུ་དེ་སྤེའུ་ཕྲུག་དང་\voc{འདྲ་}{to be like}བ་ཡིན་ནོ། རྗེས་བུ་ལ་སྨྲས་པ། \voc{ཁྱེད་}{you}ཆུང་ངུ་ཡིན་ཡང་བྱམས་པ་དང་ལྡན་པར་གྱིས་ཤིག སེམས་ཅན་ལ་ཉེས་པ་མ་བྱེད་ཅིག དེ་བཞིན་དུ་ཆོས་སྟོན་ཅིག་ཅེས་སོ། བུས་ཕ་ལ་བལྟས་ཏེ། བདག་གིས་ཁྱེད་ཀྱི་\voc{བཀའ་}{order}བཞིན་དུ་བྱའོ་ཞེས་སོ།

དེའི་ཚེ་\voc{བལ་པོ་}{Nepalese}ཡུལ་ནས་ཁྱིམ་བདག་མང་པོ་འོངས། བལ་པོ་དེ་ཀུན་གྱིས་ལོ་རྒྱུས་མང་པོ་ཤེས། དེ་ཀུན་གྱིས་ཕོ་བྲང་གི་སྒོ་ནས་རྗེ་ལ་ཕྲིན་བྱས་ཏེ། ཁྱེད་ཀྱི་ཡུལ་ནི་ཁ་བ་ཅན་དང་བཅས་པ་ཡིན། དེས་ན་བོད་ཡུལ་\voc{ཁ་བ་ཅན་}{snowy}དུ་གྲགས་སོ། བདག་ཅག་གི་ཡུལ་དང་ཁྱེད་ཀྱི་ཡུལ་བར་ལ་རི་ཆེན་པོ་དང་ཆུ་བོ་ཡོད་དོ་ཞེས་སོ། རྗེས་དེ་ཐོས་ནས་ཡི་རངས་ཏེ། བལ་པོ་དེ་ཀུན་ལ་ཟས་དང་གོས་བཏང་ངོ་།

དེ་ནས་རྗེའི་བུ་ཆུང་ངུ་དེ་ནགས་མར་སོང་། ནགས་མའི་དྲུང་ན་\voc{སྐོང་ཆུ་ལག་}{river-arm}ཆུང་ངུ་ཞིག་ཡོད། ཆུ་དེའི་ནང་དུ་ཕྱུགས་དང་བྱ་རོག་མང་པོ་འདུག བུ་དེས་ཆུ་ལ་བལྟས་ནས། སྐོང་ཆུ་ལག་འདི་ནི་ཆུ་བོ་ཆེན་པོའི་བུ་དང་འདྲ་བ་ཡིན་ནོ་སྙམ། དེར་ཁྲ་གཅིག་འཕུར་ཏེ་བྱ་རོག་ལ་རྡོ་འཕངས་པ་བཞིན་དུ་སྒྲ་ཆེན་པོ་བྱས། བྱ་རོག་དེ་འཇིགས་ནས་ཆུའི་ནང་དུ་ལྷུང་། བུ་དེས་བྱ་རོག་ཁྱེར་ནས་ཤིང་གི་དྲུང་ན་བཞག་གོ། བྱ་རོག་ལ་ཉེས་པ་མ་བྱུང་ཡང་། དེའི་སེམས་ཅུང་ཟད་མི་བདེའོ།

བུ་དེ་ཕོ་བྲང་དུ་ཕྱིར་འོངས་ནས་རྗེ་ལ་དེ་སྐད་སྨྲས། རྗེས་དེ་ཐོས་ཏེ། བྱ་རོག་ཆུང་ངུ་ཡང་ཚེ་སྲོག་དང་ལྡན་པ་ཡིན། གང་ལ་ཚེ་སྲོག་ཡོད་པ་དེ་ལ་སྡུག་བསྔལ་མ་གཏོང་ཞིག ཆུ་བོའི་ནང་དུ་ལྷུང་ཡང་། བྱམས་པ་དང་ལྡན་པས་ཕྱིར་འབྱིན་པར་བྱའོ་ཞེས་སོ། བུས་རྗེའི་བཀའ་ཐོས་ནས་སྤྲོ་བ་ཆེན་པོ་སྐྱེས་ཏེ། བདག་གིས་སེམས་ཅན་ཐམས་ཅད་ལ་བྱམས་པར་བྱའོ་སྙམ་མོ།

དེའི་བར་ལ་བལ་པོ་གཅིག་གིས་ཕོ་བྲང་གི་ནང་དུ་ཡི་གེ་ཞིག་བསྐུར། ཡི་གེ་དེར་བོད་ཀྱི་\voc{ལོ་རྒྱུས་}{account}ཐུང་ངུ་ཞིག་ཡོད། དེའི་ནང་དུ་\voc{གཉའ་ཁྲི་བཙན་པོ་}{Gnya-khri}ཞེས་བྱ་བའི་རྗེ་བྱོན་ནོ་ཞེས་ཡོད། བུ་ཆུང་ངུས་ཡི་གེ་དེ་བཀླགས་ནས། སྔོན་གྱི་རྗེ་དང་ད་ལྟར་གྱི་རྗེ་ཐམས་ཅད་ཀྱིས་ཡུལ་སྐྱོང་བར་བྱ། རྗེ་ཡིན་ཡང་སེམས་ཅན་ལ་ཉེས་པ་མ་གཏོང་ཞིག་སྙམ་མོ། རྗེས་བུ་ལ་སྨྲས་པ། ལོ་རྒྱུས་ཤེས་ན་སྔོན་གྱི་ལས་ཀྱང་མཁྱེན་ནོ། འོན་ཀྱང་སྔོན་གྱི་ལས་ཤེས་ཀྱང་། ད་ལྟར་བྱམས་པ་མེད་ན་ཆོས་མེད་དོ་ཞེས་སོ།

དེ་ནས་དེ་བཞིན་གཤེགས་པ་ཕོ་བྲང་དེར་བྱོན་ནོ། རྗེ་དང་བུ་སྨད་དང་བལ་པོ་ཐམས་ཅད་ཀྱིས་\voc{དེ་བཞིན་གཤེགས་པ་}{Tathāgata}ལ་ཕྱག་བྱས། དེ་བཞིན་གཤེགས་པ་ཤིང་གི་དྲུང་ན་བཞུགས་ཏེ། འཁོར་བ་ན་སྐྱེ་བ་དང་འཇིག་པ་ཡོད། སྐྱེས་པ་ཐམས་ཅད་འགྲོ་བར་འགྱུར། རྒྱལ་པོ་ཡང་འགྲོ། སེམས་ཅན་ཆུང་ངུ་ཡང་འགྲོ། དེས་ན་བྱམས་པ་དང་ལྡན་པར་གྱིས་ཤིག ཉེས་པ་མ་བྱེད་ཅིག་ཅེས་ཆོས་བསྟན་ཏོ།

རྗེའི་བུ་དེས་དེ་ཐོས་ཏེ་སྤྲོ་བ་དང་བྱམས་པ་སྐྱེས། བུ་དེས་སེ་འུ་བླངས་ནས་དེ་བཞིན་གཤེགས་པ་ལ་ཕུལ་ལོ། དེ་བཞིན་གཤེགས་པས་སེ་འུ་དེ་མ་བཟས་ཀྱང་། སེམས་ཅན་ལྟོགས་པ་ལ་གཏོང་བར་བཀའ་བསྩལ། བུས་བཀའ་བཞིན་དུ་སེ་འུ་དེ་བཏང་ངོ་། སེམས་ཅན་ལྟོགས་པས་སེ་འུ་དེ་\voc[ཟ་]{ཟོས་}{to eat}ནས་ཚིམ་པར་གྱུར། བུ་དེས་དེ་མཐོང་ནས། དངོས་པོ་ཅུང་ཟད་བཏང་ཡང་སྤྲོ་བ་ཆེན་པོ་སྐྱེས་སོ་སྙམ་མོ།

ཕོ་བྲང་གི་མཐའ་ན་དུར་ཞིག་ཡོད། \voc{དུར་}{grave}དེའི་དྲུང་ན་རྡོ་དང་རུས་བུ་མང་པོ་འདུག བུ་ཆུང་ངུ་དེ་དུར་དེར་སོང་ནས་རུས་བུ་ལ་བལྟས། དེར་སྔོན་གྱི་སེམས་ཅན་གྱི་ལུས་ས་ལ་ཐིམ་པ་མཐོང་ནས། ལུས་ཐམས་ཅད་\voc{ཐིམ་}{to disappear}པར་འགྱུར། སྐྱེ་བ་དང་འཇིག་པ་ནི་འཁོར་བའི་ཆོས་ཡིན་ནོ་སྙམ་མོ། བུའི་སེམས་ཅུང་ཟད་མི་དགའ་ཡང་། དེའི་བར་ལ་དེ་བཞིན་གཤེགས་པའི་ཚིག་དྲན་ཏེ་སླར་སྤྲོ་བ་སྐྱེས་སོ།

དེ་ནས་དགྲ་ཞིག་ཕོ་བྲང་གི་སྒོར་འོངས་ཏེ་རྗེ་ལ་གཤེས། དགྲ་དེས། ཁྱེད་ཀྱི་བཀའ་ངན་པ་ཡིན། ཁྱེད་ཀྱིས་ནོར་གཞན་ལ་བཏང་། ཁྱེད་རང་ཉིད་ལ་ནོར་མེད་དོ་ཞེས་\voc{གཤེ་}{to inveigh}འོ། རྗེས་དགྲ་ལ་མ་གཤེས། རྗེས་སྨྲས་པ། ནོར་བཏང་ཡང་བྱམས་པ་མི་ཟད། ཟས་བཏང་ཡང་ཆོས་མི་ཟད། ཁྱོད་ཀྱིས་གཤེས་ཀྱང་བདག་གི་སེམས་མི་འཇིག ཆོས་བཞིན་དུ་བྱེད་དོ་ཞེས་སོ། དགྲ་དེས་རྗེའི་ཚིག་ཐོས་ནས་སེམས་འཇམ་པར་གྱུར།

དེའི་ཚེ་སྒྲ་ཆེན་པོའི་ཡུལ་དུ་ཆོས་དར་རོ། ཁ་བ་ཅན་གྱི་རི་ལས་ཆུ་བོ་བྱུང་། ཆུ་བོ་དེ་ནང་དུ་སྐོང་ཆུ་ལག་མང་པོ་ཡོད། ཆུ་བོ་དང་སྐོང་ཆུ་ལག་བར་ལ་ཤིང་ཐོག་དང་སེ་འུ་མང་པོ་སྐྱེས། མི་དང་བྱ་དང་ཕྱུགས་ཐམས་ཅད་ཀྱིས་ཟས་ཟོས་ཀྱང་། ཉེས་པ་མ་བྱས། རྗེའི་བུས་དེ་མཐོང་ནས། ཆོས་བཞིན་དུ་བྱས་ན། སྤྲོ་བ་དང་བྱམས་པ་རྒྱས་སོ་སྙམ་མོ།
\end{readersection}

\chapter{ཁ་སྣོན་ཀློག་ཚན་}
\end{document}
